\documentclass[12pt,a4paper]{article}
\IfFileExists{pt-common.sty}{\usepackage{pt-common}}{\usepackage{../shared/pt-common}}
\usepackage[margin=2.5cm]{geometry}
\usepackage{setspace}
\usepackage{tikz-3dplot}
\onehalfspacing

\title{\textbf{The Theory of Persistence~I:\\The Sieve as a Dynamical System}}
\author{Yan Senez\\\texttt{yan.senez@protonmail.com}}
\date{April 2026}

\begin{document}
\maketitle

\begin{abstract}
This is the first of five articles presenting the mathematical
foundations of the Theory of Persistence (PT).  The single input is
the symmetry parameter $s = 1/2$, forced by the forbidden
self-transitions of the mod-$3$ residue sequence among $6$-rough
integers (Theorem~T1): the transfer matrix
$T_3 = \mathrm{antidiag}(1,1)$ has a unique stationary distribution
at $(1/2, 1/2)$, with no free parameter.  This is not an assumption:
$s = 1/2$ is the unique fixed point of a $2 \times 2$ doubly
stochastic antidiagonal matrix (linear algebra), and the Dirichlet
equidistribution theorem confirms the empirical convergence.

From this starting point we establish four groups of results.
\textbf{(i)}~The sieve as a dynamical field (Section~\ref{sec:sieve}):
the gap sequence $\{g_n\}$ is the unique degree of freedom;
T1 forces $T_3$ and $s = 1/2$ (unique stationary distribution of the
antidiagonal matrix).
\textbf{(ii)}~Uniqueness of Eratosthenes (Section~\ref{sec:uniqueness}):
four naturality theorems (N1--N4) prove that Eratosthenes is the unique
self-consistent, minimal, first-dynamical sieve; Theorem~T6 provides
three methodologically distinct proofs---field-theoretic (T6a), axiomatic via four
ring-congruence axioms C1--C4 (T6b), and information-geometric via
Shore--Johnson and \v{C}encov uniqueness (T6c).
\textbf{(iii)}~Conservation laws (Section~\ref{sec:conservation}):
the maximum-entropy gap distribution is geometric (L0, by strict
concavity); the spectral conservation exponent satisfies
$\alpha_{\rm cons} = s^2 = 1/4$ (Theorem~T2); the vertex--edge
bifurcation produces exactly two canonical parameters
$\qrel = 13/15$ and $\qtherm = e^{-1/15}$ at the fixed point
$\mustar = 15$ (proved in Article~III~\cite{companion_M3}).
\textbf{(iv)}~The Gap Fundamental Identity (GFT) (Section~\ref{sec:gft}):
the algebraic identity (valid for any distribution on $m$ classes)
$\log_2 m = \DKL + H$ decomposes the
total information capacity into persistent structure and irreducible
noise; the Arrow of Time $dH/d\DKL = -1$ is an identity, not an
inequality; the Ruelle identification connects the GFT to the
variational principle of thermodynamic formalism with free energy
$F_R = 0$.

All results use finite combinatorics, linear algebra, and the
Perron--Frobenius theorem.  No physics is invoked; no parameter is
fitted.  The four subsequent articles develop information geometry
and holonomy (M2), convergence and the fixed point (M3), extended
mathematical structures (M4), and the bridge from arithmetic to
physics (M5).
\end{abstract}

\tableofcontents
\newpage

% ============================================================
%  Introduction
% ============================================================
\section{Introduction}
\label{sec:intro}

This article presents the first layer of the mathematical foundations
of the Theory of Persistence (PT), a framework that studies the
information-theoretic and dynamical structure of the prime gap sequence
through the Eratosthenes sieve.  The central observation is that the
sieve imposes exact constraints on residue transitions---most
fundamentally, the absolute prohibition of self-transitions modulo~3
(Theorem~T1)---from which a rich mathematical structure unfolds with
no free parameter.

The development is organised as a five-article series:
\begin{description}[nosep]
\item[M1 (this article):] The sieve as a dynamical field, uniqueness
  of Eratosthenes, conservation laws, and the Gap Fundamental Theorem.
\item[M2:] Information geometry, holonomy, and the internal angle scale.
\item[M3:] Convergence, exact recurrences, and the fixed point
  $\mustar = 15$.
\item[M4:] Extended mathematical structures---sieve algebra, complex
  mechanics, categorical framework, predictive mathematics.
\item[M5:] The bridge from arithmetic to physics---Pontryagin duality,
  Buchstab contraction, Osterwalder--Schrader reconstruction.
\end{description}

The present article establishes the arithmetic core.  We begin by
identifying the gap sequence $\{g_n = p_{n+1} - p_n\}$ as the unique
dynamical field (Section~\ref{sec:sieve}).  The forbidden-transition
law T1 forces the mod-$3$ transfer matrix to be purely antidiagonal,
with unique stationary distribution $s = 1/2$.  Three independent routes to
$T_3$ are given (gap arithmetic, $\mathbb{Z}/6\mathbb{Z}$ involution,
spectral constraint).

Section~\ref{sec:uniqueness} proves that the Eratosthenes sieve is
the \emph{unique} sieve compatible with the multiplicative structure
of~$\mathbb{N}$.  Four theorems of increasing strength (N1--N4)
establish algebraic uniqueness, self-consistency, structural minimality,
and the canonical ordering that forces $p = 3$ as the first cascade
level.  Theorem~T6 then provides three methodologically distinct proofs of sieve
uniqueness: from field structure (T6a), from four ring-congruence
axioms (T6b), and from canonical information geometry (T6c).

Section~\ref{sec:conservation} derives the conservation laws.  The
unique maximum-entropy gap distribution is geometric (L0); the spectral
conservation exponent $\alpha_{\rm cons} = s^2 = 1/4$ controls
convergence rates (T2); the vertex--edge bifurcation at $\mustar = 15$
(proved in Article~III~\cite{companion_M3})
produces exactly two canonical deformation parameters.

Section~\ref{sec:gft} establishes the Gap Fundamental Identity (GFT):
$\log_2 m = \DKL(P\|U_m) + H(P)$, Shannon's decomposition applied to
prime residue distributions.  The mathematical content is an algebraic
identity valid for \emph{any} probability distribution on $m$ classes;
the novelty is the \emph{application} to sieve dynamics, not the
algebra.  The Ruelle equivalence, Arrow of Time, and Bekenstein bound
follow as immediate corollaries.

All proofs use only finite combinatorics, linear algebra, and the
Perron--Frobenius theorem.  No physics is invoked anywhere.  The
numbered theorems are: T1 (forbidden transitions), T3 (antidiagonal
transfer), T6 (sieve uniqueness, three proofs), L0 (max-entropy),
T2 (spectral conservation), GFT (information identity), and the
naturality chain N1--N4.


% ============================================================
%  S2. The Sieve as a Dynamical Field
% ============================================================
\section{The Sieve as a Dynamical Field}
\label{sec:sieve}

\epigraph{God made the integers; all else is the work of man.}{Leopold Kronecker}

% ---- Status Box (R1) ----------------------------------------
\begin{statusbox}
\textbf{GOAL:}
  Establish the gap sequence $\{g_n\}$ as the unique dynamical field
  of the Theory of Persistence and derive the exact mod-3 transition law T1.

\textbf{INPUTS:}
  None. This is the entry point.

\textbf{MAIN RESULT:}
  T1 (Theorem~\ref{thm:T1}): self-transitions mod~3 are absolutely forbidden.
  $T_3 = \mathrm{antidiag}(1,1)$ (Theorem~\ref{thm:T3}).
  $s = 1/2$ (unique stationary distribution of $T_3$; Theorem~\ref{thm:T3}).

\textbf{STATUS:}
  \THMTAG~(T1, T3, $s = 1/2$, Unique DOF, T1-3gram)

\textbf{NOT PROVED:}
  Convergence $\alpha_k \to 1/2$ (that is T5, the companion article~\cite{companion_M3}).
  Uniqueness of the Eratosthenes sieve (\cref{sec:uniqueness}).
  Any physical claim.
\end{statusbox}

\begin{remark}[What PT does \emph{not} claim]
\label{rem:ch1_limits}
Before proceeding, we state explicitly what lies \emph{outside}
the scope of this work:
\begin{enumerate}[nosep]
\item PT does not claim to be a ``Theory of Everything.''
  The bridge to physical observables is developed in
  Article~V~\cite{companion_M5}; this article derives zero physical
  observables.  PT says nothing about consciousness, ethics, or
  subjective experience.
\item PT does not prove that the physical universe \emph{is} the sieve.
  It proves that the sieve's structure matches the SM's constants.
  Whether this match is accidental, structural, or ontological is
  an open philosophical question.
\item PT does not replace quantum field theory.  It reconstructs
  the QFT framework (Hilbert space, Born rule, Feynman rules) from
  sieve data, but uses QFT's language to express physical predictions.
\item PT does not explain \emph{why} mathematics exists.  The mystery
  recedes one level: from ``why these constants?''\ to ``why these
  mathematical structures?''
\item Every result carries an explicit epistemic tag
  (\THMTAG, \IDTAG, \DERTAG, \BRIDGETAG, etc.).
  Claims labelled \BRIDGETAG\ are structural identifications, not
  theorems; claims labelled \CONDTAG\ depend on stated hypotheses.
  The reader should always check the tag before assessing a claim.
\end{enumerate}
\end{remark}

\begin{remark}[Independent convergence]
\label{rem:ch1_convergence}
The bridge between number theory and fundamental physics that PT
constructs systematically has recently attracted independent support
from several directions:
\begin{itemize}[nosep]
\item \textbf{Anomaly cancellation is arithmetic.}
  Camp, Gripaios and Nguyen~\cite{CampGripaiosNguyen2024} proved that
  abelian gauge anomaly cancellation in 2D admits solutions if and only
  if it admits solutions in $\mathbb{Q}_p$ for every prime~$p$
  (Hasse--Minkowski principle).
  Quantum consistency reduces to a number-theoretic
  constraint---exactly as PT derives it from the sieve.
\item \textbf{Zeta zeros are physical spectra.}
  Remmen~\cite{Remmen2021} constructed a scattering amplitude whose
  mass tower coincides with the Riemann zeta zeros, with the critical
  line $\Re(s) = 1/2$ playing the role of a unitarity bound.
\end{itemize}
These results were obtained without knowledge of PT and through
entirely different methods. They do not prove PT, but they establish
that its central premise---number theory dictates physics---is
gaining independent traction.
\end{remark}

% ============================================================
\subsection{Intuition: the sieve as a dynamical system}
\label{sec:ch1_intuition}

The Sieve of Eratosthenes is usually presented as an algorithm for
generating primes: strike out multiples of~2, then~3, then~5, and so
on. What remains are the primes. This is correct but incomplete. The
sieve does not merely \emph{filter}---it also \emph{constrains}. At
each stage, removing composites imposes selection rules on which
residue classes survive at every subsequent stage. These selection
rules accumulate into a dynamical system on the sequence of gaps
$g_n = p_{n+1} - p_n$ between consecutive primes.

\begin{remark}[Zero is certainty, not nothingness]
\label{rem:zero_certainty}
The sieve separates \textbf{certainty from uncertainty}.

A composite $n = kp$ satisfies $n \equiv 0 \pmod{p}$: its
$p$-cycle has \emph{closed}---returned to the origin after $k$~full
turns of $\mathbb{Z}/p\mathbb{Z}$.  We \emph{know} it is divisible
by~$p$; there is no uncertainty about its residue.  The zero residue
is not absence but \textbf{certainty}: the state of complete
phase cancellation, where the cycle is finished and the answer
is determined.

A prime $p$ is the opposite: no cycle closes below~$p$.  Its residues
modulo all smaller primes are non-zero---perpetually
\emph{open phases}, carrying unresolved uncertainty.
No modular test predicts a prime; it is \emph{irreducibly uncertain}.

The hierarchy of certainty is therefore:
\begin{center}
\begin{tabular}{cl}
\toprule
\textbf{Object} & \textbf{Status} \\
\midrule
$0$     & Absolute certainty (every cycle closed, $H = 0$) \\
$1$     & Neutral (no prime divides; no information) \\
Composites & Partial certainty (some cycles closed) \\
Primes  & Maximal uncertainty (no cycle closes) \\
\bottomrule
\end{tabular}
\end{center}
The sieve \emph{removes certainties} (closed cycles) and
\emph{retains uncertainties} (open phases).  Primes survive
precisely because they are \emph{unpredictable} from below.

This reversal---zero as certainty, primes as uncertainty---is
foundational.  It reappears as:
\begin{itemize}[nosep]
\item The GFT identity (\cref{sec:gft}): $\DKL + H = \log_2 m$,
  where $\DKL$ (structure = certainty) and $H$ (entropy = uncertainty)
  are conserved duals.
\item The binary operator $p = 2$ (\cref{rem:p2_operator}):
  the first certainty (even/odd is \emph{known}), which creates
  the info/anti-info partition.
\item The uncertainty operator $p = 3$: the first \emph{uncertainty}
  ($T_3$ antidiagonal, $s = 1/2$, unpredictable alternation),
  which creates the first spatial dimension.
\item The fixed-point equation (the companion article~\cite{companion_M3}):
  a \emph{global cycle closure}---the sieve recognises itself.
\end{itemize}
\end{remark}

\begin{remark}[Genesis: from two certainties to all of physics]
\label{rem:genesis}
The integers $0$ and $1$ are two \textbf{distinct poles of certainty}:
\begin{itemize}[nosep]
\item $0$: \emph{local} certainty---every prime divides~$0$;
  every cycle is closed; every question answered.
\item $1$: \emph{global} certainty---no prime divides~$1$;
  it is the identity, the neutral element, certain of being
  nothing-in-particular.
\end{itemize}

Between these two poles there are exactly \textbf{two states}:
divisible or not, closed or open, zero or one.  This count---$2$---is
not a postulate; it is the \emph{number of positions between the two
certainties}.  It is the first prime $p_1 = 2$: the binary operator.

Having two states that are not identical forces a
\textbf{third structure}: the transition between them.  Two states
$A, B$ that alternate create a sequence $A B A B \ldots$ whose
residues modulo~$3$ cycle through $\{1, 2, 1, 2, \ldots\}$, never
hitting the same class twice in a row (Theorem~T1 below).
The number $3 = p_2$ is the \emph{arithmetical consequence} of
having $2$ binary states.  It is the uncertainty operator: the first
level where the outcome is not predetermined by the binary structure.

The causal chain is therefore:
\begin{equation}
  \{0, 1\} \;\text{exist}
  \;\;\Longrightarrow\;\;
  2 = p_1
  \;\;\xrightarrow{\text{T1}}\;\;
  s = \tfrac{1}{2}
  \;\;\Longrightarrow\;\;
  3 = p_2
  \;\;\Longrightarrow\;\;
  5, 7 = p_3, p_4
  \;\;\Longrightarrow\;\;
  \mustar = 15
  \;\;\Longrightarrow\;\;
  \text{SM.}
  \label{eq:genesis}
\end{equation}
The \emph{true} input of the Theory of Persistence is not $s = 1/2$ but the
existence of two distinct certainties.  Everything else---the primes,
the sieve, the holonomy, the physical observables (Article~V~\cite{companion_M5})---follows from the
arithmetic of the space between $0$ and~$1$.
\end{remark}

\begin{remark}[Two directions, one consequence]
\label{rem:two_directions}
The space between $0$ and $1$ can be traversed in
\textbf{two directions}:
\[
  \underbrace{1 \to 0}_{\text{openness} \to \text{certainty}}
  \qquad\text{and}\qquad
  \underbrace{0 \to 1}_{\text{certainty} \to \text{openness}}
\]
Both see the same $2$~states between the poles; both create the
same $p_1 = 2$; both lead to the same T1, the same $s = 1/2$,
the same physics.  But they see it from \emph{opposite sides}.

These two directions \emph{are} the bifurcation
$\qrel \neq \qtherm$ (\cref{sec:conservation}):
\begin{align}
  q_{\rm rel} &= 1 - \frac{2}{\mu}
  &&\text{(\textbf{relativistic branch}: builds structure,}
  \notag\\
  &&&\text{discrete, carries the binary $p_1 = 2$)}\,, \label{eq:q_rel} \\[4pt]
  q_{\rm therm} &= e^{-1/\mu}
  &&\text{(\textbf{thermodynamic branch}: dissipates structure,}
  \notag\\
  &&&\text{continuous, Boltzmann limit)}\,. \label{eq:q_therm}
\end{align}

The relativistic branch carries the factor~$2 = p_1$ in its
definition: the binary operator creates the lattice $2\mathbb{N}$
on which the cascade operates.  It governs couplings
($\alphaEM$, PMNS), vertex amplitudes, and the info sector.
The thermodynamic branch is the continuous limit where the
lattice spacing vanishes; it governs geometry (metric, CKM),
edge propagators, and the entropy sector.

Both branches produce the \emph{same consequence}---the same
$s = 1/2$, the same primes, the same $\mustar = 15$---because
the symmetry $0 \leftrightarrow 1$ (exchange of the two
certainties) is exact.  This symmetry is \textbf{CPT}:
\begin{itemize}[nosep]
\item \textbf{C} (charge conjugation): exchange the \emph{contents}
  (info $\leftrightarrow$ anti-info).
\item \textbf{P} (parity): exchange the \emph{poles}
  ($0 \leftrightarrow 1$).
\item \textbf{T} (time reversal): exchange the \emph{directions}
  ($0{\to}1 \leftrightarrow 1{\to}0$).
\end{itemize}
CPT invariance is exact because the physics depends only on the
\emph{space between} the two certainties, not on which certainty
is called $0$ and which is called~$1$.

\medskip\noindent
\textbf{Terminology.}
The vertex parameter is denoted $\qrel$ (``relativistic''),
reflecting its role as the structure-building branch that carries
the binary operator.  The edge parameter is $\qtherm$
(``thermodynamic'').

\medskip\noindent
\textbf{Epistemic note.}
The identifications ``relativistic'' and ``thermodynamic'' are
\emph{bridge claims} (see the companion article~\cite{companion_M5}), not arithmetic
theorems.  The arithmetic content of this remark is solely that
two directions exist between $0$ and $1$, producing two distinct
values of~$q$.  The physical labels anticipate the bridge
identification developed in the companion article~\cite{companion_M5}.
\end{remark}

\begin{remark}[The sieve as addition/multiplication duality]
\label{rem:add_mult_duality}
The Sieve of Eratosthenes alternates between two operations
at each level~$p$:
\begin{itemize}[nosep]
\item \textbf{Multiplication} ($\times p$): remove all multiples of~$p$.
  This creates a \emph{periodic} pattern with period~$p$---a
  \textbf{cycle} (the wheel at level~$p$).
\item \textbf{Addition} ($+1$): advance to the next integer and
  measure the gap to the next survivor.
\end{itemize}

The interaction is constructive:
\begin{enumerate}[nosep]
\item $\times p$ creates \emph{bouclage}: the wheel for the
  primorial $p_1 \cdots p_k$ repeats with period $p_1 \cdots p_k$.
\item $+$ creates \emph{new gap values} at each level:
  new gaps are \textbf{sums} of previous gaps
  ($4 = 2{+}2$, $6 = 4{+}2$, $10 = 4{+}6$).
\item The bouclage \emph{exhausts} at $p_{k+1}^2$: beyond this
  threshold, composites of larger primes break the periodicity.
  The next prime restores it at a longer period.
\end{enumerate}

Numerically, 59 of the first 60 prime gaps are expressible as
$g_n = g_i + g_j$ or $g_n = g_i \times g_j$ of earlier gaps.
The sole exception is $g_0 = 1$ (the seed).  Multiplication
is always \emph{redundant} with addition---whenever a gap is a
product, it is also a sum---but it creates the periodicity that
addition cannot.

This duality echoes Remark~\ref{rem:zero_certainty}:
$0$ is the neutral element of $+$, $1$ is the neutral element
of $\times$, and the sieve alternates between them.
The bifurcation $\qrel \neq \qtherm$
(\cref{rem:two_directions}) is this same duality projected:
$\qrel = (\mustar{-}2)/\mustar$ (quotient, multiplicative),
$\qtherm = e^{-1/\mustar}$ (Taylor series, additive).
\end{remark}

The key insight is simple: consecutive primes cannot share a residue
class modulo~3. If $p_n \equiv r \pmod{3}$ and $p_{n+1} \equiv r
\pmod{3}$, then their difference $g_n = p_{n+1} - p_n \equiv 0
\pmod{3}$, which forces $3 \mid g_n$. For gaps of size at least~2
and primes greater than~3, this means the integer $p_n + g_n/2$
would be divisible by~3---contradicting its primality if $g_n \geq 6$,
or leading to a composite midpoint for smaller even multiples of~3.
The detailed argument is Theorem~T1 below; the full sieve
process is visualised in Figure~\ref{fig:sieve_process}.

The gap sequence $\{g_n\}$ is not merely a derived quantity. It is
the \emph{only} independent degree of freedom: given the gaps and the
initial prime $p_1 = 2$, the entire prime sequence is reconstructed;
conversely, knowing all residue classes modulo every prime recovers
the gaps. Everything in the Theory of Persistence---the transfer matrices,
the holonomy angles, the coupling constants---is a functional of
$\{g_n\}$.

\begin{remark}[PT is neither local nor global]
\label{rem:ch1_not_local}
The Chinese Remainder Theorem (CRT) factorises each transfer matrix
$T_m = T_{p_1} \otimes \cdots \otimes T_{p_k}$ into prime-by-prime
factors. This is a powerful \emph{tool}, but it does not make PT a
``local'' theory where each prime acts independently.  Several
structures are intrinsically collective:
\begin{enumerate}[nosep]
\item The H-theorem $S_{K+1} \geq S_K$: a monotonicity constraint
  on the \emph{full} distribution.
\item Spectral annihilation $r_2(0) = 0$ (the companion article~\cite{companion_M3}): a property of
  $T_3$ as a whole, not decomposable into per-prime factors.
\item The Fisher metric (the companion article~\cite{companion_M2}): a continuous Riemannian geometry
  that \emph{is} the sieve, not an approximation of it.
\item Holonomy $\sin^2(\theta_p)$ (the companion article~\cite{companion_M2}): a phase accumulated
  around $\mathbb{Z}/p\mathbb{Z}$---global coherence, not a
  pointwise quantity.
\item The automorphic invariance lemma (the companion article~\cite{companion_M3}): rigidity of
  $R_p = \{0\}$ under $\mathrm{Aut}(\mathbb{Z}/p\mathbb{Z})$.
\end{enumerate}
PT dissolves the local/global distinction just as it dissolves the
discrete/continuous one: the CRT product structure and the
collective constraints coexist as complementary descriptions of a
single arithmetic object.
\end{remark}

\begin{figure}[htbp]
\centering
\begin{tikzpicture}[>=stealth, xscale=0.60]
  % --- Row 0: All integers 2..23 ---
  \node[font=\scriptsize\bfseries, anchor=east] at (-1.4, 4.2) {Integers};
  \foreach \n [count=\i from 0] in {2,3,4,5,6,7,8,9,10,11,12,13,14,15,16,17,18,19,20,21,22,23} {
    \node[font=\scriptsize, minimum width=0.55cm] at (\i, 4.2) {\n};
  }

  % --- Row 1: Sieve — composites struck, primes circled ---
  \node[font=\scriptsize\bfseries, anchor=east] at (-1.4, 2.8) {Sieve};
  % Composites (struck out) — index = n-2
  \foreach \i/\n in {2/4, 4/6, 6/8, 7/9, 8/10, 10/12, 12/14, 13/15, 14/16, 16/18, 18/20, 19/21, 20/22} {
    \node[font=\scriptsize, text=gray!45] at (\i, 2.8) {\n};
    \draw[gray!40, thin] (\i-0.18, 2.65) -- (\i+0.18, 2.95);
  }
  % Primes (circled) — index = n-2
  \foreach \i/\n in {0/2, 1/3, 3/5, 5/7, 9/11, 11/13, 15/17, 17/19, 21/23} {
    \node[font=\scriptsize\bfseries, thmblue] (pr\n) at (\i, 2.8) {\n};
    \draw[thmblue, semithick] (\i, 2.8) circle (0.27cm);
  }

  % --- Row 2: Gaps (between primes, with arrows) ---
  \node[font=\scriptsize\bfseries, anchor=east] at (-1.4, 1.3) {Gap $g_n$};
  \foreach \xa/\xb/\g in {0/1/1, 1/3/2, 3/5/2, 5/9/4, 9/11/2, 11/15/4, 15/17/2, 17/21/4} {
    \pgfmathsetmacro{\xm}{(\xa+\xb)/2}
    \draw[condorange, <->, semithick] (\xa+0.32, 1.3) -- (\xb-0.32, 1.3);
    \node[font=\scriptsize\bfseries, condorange, fill=white, inner sep=1pt]
      at (\xm, 1.3) {\g};
  }

  % --- Row 3: Residues mod 3 ---
  \node[font=\scriptsize\bfseries, anchor=east] at (-1.4, -0.1) {$p \!\bmod 3$};
  \foreach \i/\r/\c in {0/2/predred, 1/0/valgray, 3/2/predred, 5/1/bridgepurple,
    9/2/predred, 11/1/bridgepurple, 15/2/predred, 17/1/bridgepurple, 21/2/predred} {
    \node[font=\scriptsize\bfseries, \c, draw=\c, rounded corners=2pt,
          minimum width=0.45cm, minimum height=0.38cm, inner sep=1pt]
      at (\i, -0.1) {\r};
  }

  % Annotation: alternation
  \draw[decorate, decoration={brace, amplitude=4pt, mirror}, thick]
    (3-0.35, -0.55) -- (21+0.35, -0.55);
  \node[font=\scriptsize, align=center] at (12, -1.0)
    {$p > 3$: residues alternate $1 \leftrightarrow 2$ (Theorem~T1)};

  % Vertical dashed guides from sieve row to residue row
  \foreach \i in {0, 1, 3, 5, 9, 11, 15, 17, 21} {
    \draw[gray!25, densely dotted] (\i, 2.42) -- (\i, 0.15);
  }
\end{tikzpicture}
\caption{The sieve process visualised.
  \textbf{Row~1}: integers 2--23.
  \textbf{Row~2}: composites struck out (Eratosthenes), primes circled.
  \textbf{Row~3}: gaps $g_n = p_{n+1} - p_n$.
  \textbf{Row~4}: residues mod~3---after $p=3$, they strictly alternate
  between 1 and 2, forbidden from repeating (T1).}
\label{fig:sieve_process}
\end{figure}

% ============================================================
\subsection{The gap sequence: sole degree of freedom}
\label{sec:gap_sequence}

\begin{definition}[Gap Sequence]
\label{def:gaps}
Let $2 = p_1 < p_2 < p_3 < \cdots$ be the sequence of all primes in
increasing order. The \emph{gap sequence} is
\begin{equation}
  g_n = p_{n+1} - p_n, \qquad n \geq 1.
  \label{eq:gap_def}
\end{equation}
The first values are $g_1 = 1$, $g_2 = 2$, $g_3 = 2$, $g_4 = 4$,
$g_5 = 2$, $g_6 = 4$, $g_7 = 2$, $g_8 = 4$, $g_9 = 6$, \ldots
The initial gap $g_1 = 1$ (from $p_1=2$ to $p_2=3$) is exceptional;
all subsequent gaps are even.
\end{definition}

\begin{definition}[Residue Sequence]
\label{def:residue}
For any integer modulus $m \geq 2$, the \emph{residue sequence}
modulo~$m$ is
\begin{equation}
  r_n^{(m)} = p_n \bmod m, \qquad n \geq 1.
  \label{eq:residue_def}
\end{equation}
When $m$ is clear from context, we write $r_n$ for $r_n^{(m)}$.
\end{definition}

The residue sequence is entirely determined by the gap sequence:
\begin{equation}
  r_{n+1}^{(m)} = \bigl(r_n^{(m)} + g_n\bigr) \bmod m,
  \label{eq:gauge_connection}
\end{equation}
with initial condition $r_1^{(m)} = 2 \bmod m$. This is not an
approximation; it is an exact identity.

\mTHM
\begin{theorem}[Unique Degree of Freedom]
\label{thm:unique_dof}
The gap sequence $\{g_n\}_{n \geq 1}$ is the unique dynamical field
of \PT. For any modulus $m$, all residue sequences $\{r_n^{(m)}\}$,
all transfer matrices $T_m$, all holonomy angles $\theta_p$, and all
coupling constants derived from these are measurable functionals
of~$\{g_n\}$.
\end{theorem}

\begin{proof}
The residue reconstruction~\eqref{eq:gauge_connection} is bijective:
given $\{g_n\}$ and $r_1^{(m)} = 2 \bmod m$, the entire sequence
$\{r_n^{(m)}\}$ is determined. Conversely, $g_n = r_{n+1}^{(m)} -
r_n^{(m)} \bmod m$ recovers the gap (modulo $m$). Since this holds
for all $m$ simultaneously, the gap sequence is the unique common
generator.

For non-decomposability: suppose $g_n = f_n + h_n$ with $\{f_n\}$ and
$\{h_n\}$ statistically independent. Under the mod-3 residue map, the
constraint $P[r \to r] = 0$ (Theorem~\ref{thm:T1} below) couples consecutive
pairs $(r_n, r_{n+1})$. This coupling is preserved by any function of
$g_n$, hence any sub-field $\{f_n\}$ inheriting this coupling cannot
be independent of $\{h_n\}$. The field is indecomposable.

A complete formalization via the automorphic invariance lemma is given
in the companion article~\cite{companion_M3}: conditions U1--U4 uniquely determine the gap
sequence, subsuming both reconstruction and indecomposability.
\end{proof}

\begin{remark}
The gauge-connection form~\eqref{eq:gauge_connection} is not a
metaphor. It is the \emph{defining} property of a $\mathbb{Z}/m\mathbb{Z}$-valued
lattice gauge field: the residue $r_n$ is the ``parallel transport''
of the initial state $r_1$ along the gap field $\{g_n\}$, with the
modulus $m$ selecting the gauge group. We make this precise in
the companion article~\cite{companion_M2} (holonomy).
\end{remark}

% ============================================================
\subsection{Additive coverage of prime gaps}
\label{sec:ch1_coverage}

The gap sequence has a remarkable algebraic property: it is
\emph{additively closed}.

\begin{thmbox}{Additive Generation of Wheel Gaps}
\label{thm:wheel_coverage}
For every $k \geq 1$, the set of gap values in the wheel at
level $k{+}1$ satisfies
$W_{k+1} \subseteq W_k \cup \mathrm{ADD}(W_k)$,
where $\mathrm{ADD}(G) = \{a + b : a, b \in G\}$.

\begin{proof}
When passing from level $k$ to $k{+}1$, the sieve removes
multiples of $p_{k+1}$.  Each removal \emph{merges} two adjacent
gaps $a, b \in W_k$ into a single gap $a + b$.  Gaps whose
endpoints are both retained pass unchanged.  Hence every element
of $W_{k+1}$ is either in $W_k$ or in $\mathrm{ADD}(W_k)$.
\end{proof}
\end{thmbox}

\begin{remark}[Extension to prime gaps (conjecture)]
\label{rem:coverage_conjecture}
We \emph{conjecture} that the additive coverage extends to actual
prime gaps: for every $n \geq 2$,
$g_n \in G_{n-1} \cup \mathrm{ADD}(G_{n-1})$.
This is verified computationally on all $10^7$ prime gaps below
$1.8 \times 10^8$ ($100.0000\%$ coverage; $55.5\%$ ADD-only,
$44.5\%$ also multiplicatively generated).
The passage from wheel gaps to prime gaps requires controlling
primality of endpoints, which the wheel theorem does not address.

The pool of distinct gap values grows as $|G_N| \approx 0.13\,
(\ln p_N)^{2.26}$ and satisfies $\max(G_N) / |G_N| \approx 2$
across five orders of magnitude.
The truncated density $\rho(G_N) = |G_N \cap [2, \max G_N]|
/ (\max G_N / 2) = 0.946$ for $N = 10^7$ (105 of 111 even
values present; no hole below~198).
Conditionally on Polignac's conjecture (every even integer is
a prime gap infinitely often), the Schnirelmann density of the
infinite pool $G_\infty = \bigcup_N G_N$ would equal~1,
proving the conjecture via Schnirelmann's theorem.
Unconditionally, the gap between the wheel theorem and the
full conjecture remains an open problem.
\end{remark}

% ============================================================
\subsection{Forbidden transitions: Theorem T1}
\label{sec:T1}

\mTHM
\begin{theorem}[Forbidden Transitions (T1)]
\label{thm:T1}
Among the sieve candidates at level $p = 3$ (the $6$-rough integers,
i.e., integers coprime to both~$2$ and~$3$), the self-transitions of
the mod-$3$ residue sequence are absolutely forbidden:
\begin{equation}
  P\bigl[r_n \equiv 1 \to r_{n+1} \equiv 1 \pmod{3}\bigr] = 0, \qquad
  P\bigl[r_n \equiv 2 \to r_{n+1} \equiv 2 \pmod{3}\bigr] = 0.
  \label{eq:T1}
\end{equation}
This holds for every pair of consecutive $6$-rough integers, not merely
asymptotically.
\end{theorem}

\begin{proof}
The $6$-rough integers are exactly those of the form $6k \pm 1$.
Their residues modulo~3 alternate: $6k - 1 \equiv 2$ and
$6k + 1 \equiv 1$. Consecutive candidates are separated by gaps
of 2 or 4, and in every case the residue switches between 1 and 2.
Therefore, among sieve-level-3 candidates, self-transitions are
impossible: $P[r \to r] = 0$ exactly.
\end{proof}

\begin{corollary}[T1 for consecutive primes]
\label{cor:T1_primes}
Since every prime $\geq 5$ is $6$-rough, T1 constrains the
\emph{possible} transitions between consecutive primes, but does not
forbid all same-residue pairs outright: consecutive primes can share a
mod-$3$ residue (e.g., $(61, 67)$, both $\equiv 1 \pmod{3}$) because a
composite $6$-rough integer may intervene.
\end{corollary}

\begin{remark}[PT usage of T1]
\label{rem:T0_scope}
In the PT framework, T1 fixes the transition matrix at the sieve level:
$T_3 = \mathrm{antidiag}(1,1)$, which is exact for $6$-rough integers.
The empirical $T_3$ for consecutive primes has small but nonzero diagonal
entries, suppressed by $O(1/\ln N)$.  All PT derivations
(symmetry parameter $s = 1/2$, conservation theorem, etc.)
rely on the sieve-level matrix, not on the empirical prime-level matrix.
The departure from the ideal decreases as $N$ grows (consistent with
the Elliott--Halberstam conjecture, widely believed).
\end{remark}

% ============================================================
\subsection{The exact transition matrix $T_3$}
\label{sec:T3}

\mTHM
\begin{theorem}[Antidiagonal Transfer (T3)]
\label{thm:T3}
Among the sieve candidates at level $p = 3$ (integers $\equiv \pm 1 \pmod 6$),
the transition matrix on $\{1, 2\} \pmod 3$ is exactly
\begin{equation}
  T_3 = \begin{pmatrix} 0 & 1 \\ 1 & 0 \end{pmatrix}.
  \label{eq:T3}
\end{equation}
The transition graph is shown in Figure~\ref{fig:T3_graph}.
\end{theorem}

\begin{figure}[htbp]
\centering
\begin{tikzpicture}[>=stealth, scale=1.3]
  % Two state nodes
  \node[circle, draw=thmblue, thick, fill=thmblue!10,
        minimum size=1.2cm, font=\large] (s1) at (-1.8, 0) {$1$};
  \node[circle, draw=thmblue, thick, fill=thmblue!10,
        minimum size=1.2cm, font=\large] (s2) at (1.8, 0) {$2$};
  % Transition arrows (off-diagonal = 1)
  \draw[->, thick, thmblue] (s1) to[bend left=30]
    node[above, font=\scriptsize] {$P = 1$} (s2);
  \draw[->, thick, thmblue] (s2) to[bend left=30]
    node[below, font=\scriptsize] {$P = 1$} (s1);
  % Forbidden self-loops (crossed out)
  \draw[->, gray, dashed] (s1) to[out=140, in=220, looseness=5] (s1);
  \draw[predred, very thick] (-2.7, 0.9) -- (-2.3, 0.3);
  \draw[predred, very thick] (-2.7, 0.3) -- (-2.3, 0.9);
  \draw[->, gray, dashed] (s2) to[out=40, in=-40, looseness=5] (s2);
  \draw[predred, very thick] (2.3, 0.9) -- (2.7, 0.3);
  \draw[predred, very thick] (2.3, 0.3) -- (2.7, 0.9);
  % Matrix
  \node[font=\footnotesize, align=center] at (0, -1.8)
    {$T_3 = \begin{pmatrix} 0 & 1 \\ 1 & 0 \end{pmatrix}$
     \quad (antidiagonal)};
  % Labels
  \node[font=\scriptsize, gray] at (0, 1.2)
    {residue classes $\bmod\,3$};
  \node[font=\tiny, predred] at (-2.5, 0.0) {T1};
  \node[font=\tiny, predred] at (2.5, 0.0) {T1};
\end{tikzpicture}
\caption{Transition graph of the mod-3 sieve dynamics.
  Sieve candidates alternate between residue classes $1$ and $2$
  with probability~1. Self-transitions are absolutely forbidden (T1,
  red crosses). The transfer matrix $T_3$ is purely antidiagonal---
  the simplest non-trivial dynamical law.}
\label{fig:T3_graph}
\end{figure}

\begin{proof}
The sieve candidates at level~3 are the integers coprime to $2 \times 3 = 6$,
i.e., integers of the form $6k \pm 1$ for integer $k \geq 1$. Their residues
modulo~3 are: $6k + 1 \equiv 1 \pmod 3$ and $6k - 1 \equiv 2 \pmod 3$.
Consecutive such candidates are $(6k-1, 6k+1)$ (gap~2) or $(6k+1, 6(k+1)-1) = (6k+1, 6k+5)$
(gap~4). In both cases, residues alternate: $1 \to 2$ or $2 \to 1$. No self-transition
ever occurs among sieve-level-3 candidates. Stochasticity forces off-diagonal entries to~1.
\end{proof}

\begin{remark}[Three independent routes to $T_3 = \mathrm{antidiag}(1,1)$]
\label{rem:T3_three_routes}
The antidiagonal form of $T_3$ can be established by three disjoint arguments,
each using different mathematical primitives.

\smallskip\noindent
\textbf{Route 1: Gap arithmetic (above).}
Enumerate the candidates $6k \pm 1$; the only possible gaps are $2$ and $4$,
both of which swap the residue class modulo~3.

\smallskip\noindent
\textbf{Route 2: $\mathbb{Z}/6\mathbb{Z}$ involution.}
The survivors modulo~$6$ are $\{1, 5\}$. The successor map $\sigma$ acts on
this two-element set and is therefore the unique non-trivial involution:
$\sigma(1) = 5,\; \sigma(5) = 1$. Reducing modulo~3:
$1 \mapsto 2,\; 2 \mapsto 1$, so
$T_3 = \bigl(\begin{smallmatrix} 0 & 1 \\ 1 & 0 \end{smallmatrix}\bigr)$.
This proof uses only $|\mathrm{surv}_6| = 2$ and the involution structure,
with no gap computation.

\smallskip\noindent
\textbf{Route 3: Spectral constraint.}
$T_3$ is a $2 \times 2$ doubly stochastic matrix (by the exchange symmetry
$1 \leftrightarrow 2$) with $\mathrm{tr}(T_3) = 0$ (Theorem~T1:
self-transitions forbidden). The eigenvalues are $\{+1, -1\}$.
The unique doubly stochastic matrix with these eigenvalues is
$\mathrm{antidiag}(1,1)$.
\end{remark}

\mCL
\begin{thmbox}{Symmetry parameter: $s = 1/2$}
\label{thm:s_half}
The antidiagonal matrix $T_3 = \mathrm{antidiag}(1,1)$ (Theorem~T3)
has a unique stationary distribution:
\begin{equation}
  s = \pi_1 = \pi_2 = \frac{1}{2}.
  \label{eq:s_half}
\end{equation}

\begin{proof}
The eigenvector equation $\pi\,T_3 = \pi$ gives
$(\pi_1, \pi_2) \begin{psmallmatrix} 0 & 1 \\ 1 & 0 \end{psmallmatrix}
= (\pi_2, \pi_1) = (\pi_1, \pi_2)$,
hence $\pi_1 = \pi_2$.  Normalisation $\pi_1 + \pi_2 = 1$ forces
$s = 1/2$.  Uniqueness: the eigenvalue $\lambda_0 = 1$ of the
$2 \times 2$ doubly stochastic matrix $T_3$ has algebraic
multiplicity~1 (the other eigenvalue is $\lambda_1 = -1$).
By Perron--Frobenius, the stationary distribution is unique.
\end{proof}
\end{thmbox}

\begin{remark}[Epistemics of $s = 1/2$]
\label{rem:s_half}
The value $s = 1/2$ is a \textbf{theorem}, not an assumption.  It is
the unique fixed point of the antidiagonal transfer matrix $T_3$,
proved above by linear algebra alone.  The Dirichlet
equidistribution theorem independently confirms that empirical
frequencies converge to this value: among primes $> 3$, exactly half
lie in each of the two non-trivial residue classes modulo~3.
The \emph{rate} of convergence ($\alpha_k \to 1/2$) is quantified in
Article~III~\cite{companion_M3} via spectral annihilation.
\end{remark}

% ============================================================
\subsection{The gauge connection}
\label{sec:gauge}

Equation~\eqref{eq:gauge_connection} generalizes to arbitrary square-free
moduli via the Chinese Remainder Theorem (CRT).

\begin{definition}[Square-free modulus]
A positive integer $m = q_1 \cdots q_k$ is \emph{square-free} if
$q_1 < q_2 < \cdots < q_k$ are distinct primes.
\end{definition}

\mTHM
\begin{theorem}[Gauge Connection and CRT Factorization]
\label{thm:gauge}
Let $m = q_1 \cdots q_k$ be square-free. The residue recurrence
\[
  r_{n+1}^{(m)} = \bigl(r_n^{(m)} + g_n\bigr) \bmod m
\]
decomposes via the isomorphism $\mathbb{Z}/m\mathbb{Z} \cong
\prod_{i=1}^k \mathbb{Z}/q_i\mathbb{Z}$ (CRT) into $k$ independent
components, each evolving as $r_{n+1}^{(q_i)} = (r_n^{(q_i)} + g_n)
\bmod q_i$. The components at distinct primes $q_i \neq q_j$ are
statistically independent (to leading order in $1/\ln N$).
\end{theorem}

\begin{proof}
The CRT isomorphism is a ring isomorphism, hence exact. It maps
$r \bmod m$ to the vector $(r \bmod q_1, \ldots, r \bmod q_k)$.
The recurrence $r_{n+1} = r_n + g_n \pmod m$ maps to
$r_{n+1}^{(q_i)} = r_n^{(q_i)} + g_n \pmod{q_i}$ for each $i$.
Independence at leading order follows from the prime number theorem
for arithmetic progressions (PNT-AP): prime residues modulo $q_i$ and
$q_j$ (for distinct primes $q_i, q_j$) are asymptotically equidistributed
independently. The correction is $O(1/\ln N)$ by the PNT-AP error term.
\end{proof}

% ============================================================
\subsection{DPI monotonicity}
\label{sec:dpi}

\mID
\begin{idbox}[Data Processing Inequality]
\label{id:dpi}
If $m_1 \mid m_2$ (i.e., $m_1$ divides $m_2$), then
\begin{equation}
  \DKL(P_{m_1} \| U_{m_1}) \leq \DKL(P_{m_2} \| U_{m_2}),
  \label{eq:dpi}
\end{equation}
where $P_m$ is the prime-residue distribution at modulus~$m$ and
$U_m$ is the uniform distribution on $(\mathbb{Z}/m\mathbb{Z})^\times$.

\begin{proof}
The natural projection $\pi_{m_2 \to m_1} : \mathbb{Z}/m_2\mathbb{Z}
\to \mathbb{Z}/m_1\mathbb{Z}$ (reduction modulo $m_1$) is a surjective
ring homomorphism, hence a Markov kernel. The data processing inequality
for KL divergence states: for any stochastic map $K$,
$\DKL(K(P) \| K(Q)) \leq \DKL(P \| Q)$. Applying to $K = \pi_{m_2 \to m_1}$
with $P = P_{m_2}$ and $Q = U_{m_2}$, and noting $\pi_{m_2 \to m_1}(U_{m_2})
= U_{m_1}$ (uniform maps to uniform), we get $\DKL(P_{m_1} \| U_{m_1})
\leq \DKL(P_{m_2} \| U_{m_2})$.
\end{proof}
\end{idbox}

DPI monotonicity says: more sieving = more structure = higher divergence
from uniformity. The sieve levels are ordered by $\DKL$.

% ============================================================
\subsection{T1-3gram: three-step forbidden patterns}
\label{sec:T1_3gram}

The forbidden-transition law T1 implies constraints on three-step
patterns $(r_n, r_{n+1}, r_{n+2})$ of mod-3 residues.

\mTHM
\begin{theorem}[Trigram Validation (T1-3gram)]
\label{thm:T1_3gram}
Among sieve-level-3 candidates, the following three-step patterns are
absolutely forbidden:
\begin{equation}
  n_3(1,0,1) = 0, \qquad n_3(2,0,2) = 0.
  \label{eq:T1_3gram}
\end{equation}
More precisely: of the $27 = 3^3$ possible mod-3 triples, exactly
$12$ are forbidden:
\begin{itemize}
\item $10$ by T1 (auto-transitions within the triple), and
\item $2$ additional by the parity-alternating constraint.
\end{itemize}
The $15$ remaining triples are allowed. Time-reversal symmetry holds
exactly at the 3-gram level; $1 \leftrightarrow 2$ symmetry is broken
at the 3-gram level.
\end{theorem}

\begin{proof}
A triple $(r_n, r_{n+1}, r_{n+2})$ is forbidden if either
$(r_n, r_{n+1})$ or $(r_{n+1}, r_{n+2})$ is a forbidden pair
(T1: self-transition).

\emph{Step 1: T1-forbidden triples.}
Among the $3^2 = 9$ pairs, $2$ are forbidden ($1\to1$ and $2\to2$).
\begin{itemize}[nosep]
\item Triples with $(r_n, r_{n+1}) \in \{(1,1),(2,2)\}$:
  $2 \times 3 = 6$ (any third element).
\item Triples with $(r_{n+1}, r_{n+2}) \in \{(1,1),(2,2)\}$:
  $3 \times 2 = 6$.
\item Overlap (both pairs forbidden): $(1,1,1)$ and $(2,2,2)$: $2$.
\end{itemize}
By inclusion-exclusion: $6 + 6 - 2 = 10$ triples forbidden by T1.

\emph{Step 2: parity-forbidden triples.}
The alternating parity of sieve-level-3 candidates forbids patterns
that pass through residue~0 (excluded from the sieve). The two
additional forbidden triples are $(1,0,1)$ and $(2,0,2)$.

\emph{Conclusion.}
$10 + 2 = 12$ forbidden, $27 - 12 = 15$ allowed.
\end{proof}

% ============================================================
\subsection{Summary and connections}
\label{sec:ch1_summary}

\begin{ledgerbox}
Hard core:
  T1 (exact at sieve level), T3 = antidiag(1,1) (exact), Unique DOF,
  Gauge Connection (CRT), DPI monotonicity, T1-3gram (12/27 forbidden).

PT-Closure:
  s = 1/2 depends on stationarity (T4, 10/10 via spectral annihilation,
  see the companion article~\cite{companion_M3}) and exchange symmetry (Dirichlet equidistribution,
  not proved here).

Bridge claim:
  None in this section. No physics invoked.
\end{ledgerbox}

\medskip

The gap sequence $\{g_n\}$ is the entry point of the Theory of Persistence.
It is the sole degree of freedom (Theorem~\ref{thm:unique_dof}),
it encodes a $\mathbb{Z}/m\mathbb{Z}$ gauge connection
(Theorem~\ref{thm:gauge}), and its mod-3 behavior is constrained
by the exact alternation law (Theorems~T1 and~T3).

The next section asks: is the Eratosthenes sieve the \emph{only}
sieve compatible with these properties? The answer (T6) is yes.

\medskip\noindent\textbf{Forward connections.}
\begin{itemize}
\item \cref{sec:uniqueness}: T6 uses the CRT structure proved here.
\item \cref{sec:conservation}: Conservation laws use the transfer matrix $T_3$.
\item The companion article~\cite{companion_M2}: Holonomy angles $\theta_p$ are derived from the same gauge structure.
\item The companion article~\cite{companion_M3}: Convergence $\alpha_k \to 1/2$ uses T1 and the gauge recurrence.
\end{itemize}


% ============================================================
%  S3. Uniqueness of Eratosthenes
% ============================================================
\section{Why Eratosthenes Is the Natural Sieve}
\label{sec:uniqueness}

\epigraph{The laws of nature are but the mathematical thoughts of God.}{Euclid}

% ---- Status Box (R1) ----------------------------------------
\begin{statusbox}
\textbf{GOAL:}
  Prove that the Sieve of Eratosthenes is the \emph{unique} sieve
  compatible with the multiplicative structure of~$\mathbb{N}$
  (naturality theorems N1--N4) and with a formal axiom system
  on ring congruences (T6, C1--C4).

\textbf{INPUTS:}
  \cref{sec:sieve} (gauge connection, CRT, DPI).

\textbf{MAIN RESULTS:}
  N1--N4: Eratosthenes is the unique self-consistent, minimal,
  first-dynamical sieve.
  T6a: $R_p = \{0\}$ from field structure.
  T6b: C1--C4 $\Rightarrow$ Eratosthenes unique
  (independence 4/4 formal).
  T6c: $\DKL$ unique (Shore--Johnson), Fisher unique
  (\v Cencov).
  Overall: T6 complete.

\textbf{STATUS:}
  \THMTAG~(N1, N2, N3, N4, T6a, T6b, T6c)

\textbf{NOT PROVED:}
  Any physical interpretation of the uniqueness.
  The bridge from $N_c = 3$ to color (the companion article~\cite{companion_M5}).
\end{statusbox}

% ============================================================
\subsection{Intuition: forcing the sieve from arithmetic alone}
\label{sec:ch2_intuition}

The Sieve of Eratosthenes is usually thought of as a clever algorithm.
But it is more: it is the \emph{only} algorithm compatible with the
multiplicative structure of the integers. This section proves four
increasingly strong versions of this claim (N1--N4) and then establishes
the formal ring-congruence uniqueness (T6b).

The logical chain is:
\begin{enumerate}
\item \textbf{N1}: Primes are the unique atoms of $(\mathbb{N}, \times)$
  (from the Fundamental Theorem of Arithmetic).
\item \textbf{N2}: $\mathbb{P}$ is the unique self-consistent sieve (fixed point).
\item \textbf{N3}: Eratosthenes is structurally minimal (simplest monoid,
  weakest operation, canonical ordering).
\item \textbf{N4}: $p = 3$ is forced as the first \emph{cascade} level
  ($p = 2$ is the info/anti-info operator; \cref{rem:p2_operator});
  this gives $s = 1/2$ as the first dynamical output.
\end{enumerate}
Finally, T6b derives $E_d = [0] = d\mathbb{Z}$ as a theorem
from four axioms C1--C4, without assuming the modular form of the sieve.

% ============================================================
\subsection{Setup: sieves on monoids}
\label{sec:ch2_setup}

\begin{definition}[Multiplicative sieve]
\label{def:mult_sieve}
A \emph{multiplicative sieve} is a pair $(G, \mathcal{S})$ where
$G \subseteq \mathbb{N}_{\geq 2}$ is a set of generators and
\[
  \mathcal{S}(G) = \{n \in \mathbb{N}_{\geq 2} :
    n \text{ is not divisible by any } g \in G \text{ with } g < n\}.
\]
\end{definition}

\begin{definition}[Self-consistent sieve]
\label{def:self_consistent}
A multiplicative sieve $(G, \mathcal{S})$ is \emph{self-consistent} if
$\mathcal{S}(G) = G$.
\end{definition}

% ============================================================
\subsection{N1: Algebraic uniqueness of primes}
\label{sec:N1}

\mTHM
\begin{theorem}[Algebraic Uniqueness (N1 --- classical)]
\label{thm:N1}
The prime numbers $\mathbb{P}$ are the unique atoms of the multiplicative
monoid $(\mathbb{N}_{\geq 1}, \times)$. The Eratosthenes sieve is the
unique constructive procedure identifying these atoms by successive
elimination.
This is a classical consequence of the Fundamental Theorem of Arithmetic.
We state it for completeness.

\begin{proof}
\textbf{Part 1 (Atoms).} By the Fundamental Theorem of Arithmetic
(FTA: every $n \geq 2$ has a unique factorization into primes), an
element $p \geq 2$ is an atom if and only if it is prime. FTA is an
external import (classical theorem, Hardy--Wright~\cite{HardyWright2008}, \S1.10).

\textbf{Part 2 (Procedure).} To identify atoms of $(\mathbb{N}_{\geq 2},|)$
constructively, one must test: for each $n$, is there $a, b \geq 2$ with
$n = ab$? This requires testing divisibility by all atoms smaller than $n$.
Since the atoms are initially unknown, this bootstraps: process elements
in increasing order (well-ordering of $\mathbb{N}$), classify each as atom
or composite based on already-classified atoms. This is exactly the
Eratosthenes sieve. Any reordering either (a) produces the same result
by FTA uniqueness, or (b) requires additional structure beyond
$(\mathbb{N}, \times)$.
\end{proof}
\end{theorem}

\begin{remark}
FTA fails in $\mathbb{Z}[\sqrt{-5}]$ (where $6 = 2 \cdot 3 = (1+\sqrt{-5})(1-\sqrt{-5})$).
Eratosthenes is therefore tied to $\mathbb{N}$ as the simplest UFD.
\end{remark}

% ============================================================
\subsection{N2: Self-consistency (unique fixed point)}
\label{sec:N2}

\mTHM
\begin{theorem}[Self-Consistency (N2 --- classical)]
\label{thm:N2}
$\mathbb{P}$ is the unique self-consistent multiplicative sieve on
$\mathbb{N}_{\geq 2}$.
This is a classical consequence of the definition of primality.
We state it for completeness.

\begin{proof}
\textbf{Existence.} $n \in \mathcal{S}(\mathbb{P})$ iff $n$ has no prime
factor $p < n$. If $n$ is composite, $n = ab$ with $a \leq \sqrt{n}$;
then $a$ has a prime factor $p \leq a < n$, so $p \mid n$ and $n
\notin \mathcal{S}(\mathbb{P})$. If $n$ is prime, no prime $p < n$
divides $n$. Hence $\mathcal{S}(\mathbb{P}) = \mathbb{P}$.

\textbf{Uniqueness.} Let $G \neq \mathbb{P}$ satisfy $\mathcal{S}(G) = G$.

\textit{Case 1: $G \subsetneq \mathbb{P}$.} Some prime $p \notin G$.
Since $p$ is prime, no element of $G$ properly divides $p$; so $p \in
\mathcal{S}(G) = G$. Contradiction.

\textit{Case 2: $G \supsetneq \mathbb{P}$.} Some composite $c \in G$.
Write $c = p \cdot m$ with $p$ prime. Then $p \in \mathbb{P} \subseteq G$,
and $p < c$ divides $c$. So $c \notin \mathcal{S}(G) = G$. Contradiction.

\textit{Case 3: $G \not\subseteq \mathbb{P}$ and $\mathbb{P} \not\subseteq G$.}
Contains composites (Case~2) or misses primes (Case~1). Either leads to
contradiction. Hence $G = \mathbb{P}$.
\end{proof}
\end{theorem}

% ============================================================
\subsection{N3: Structural minimality}
\label{sec:N3}

We compare sieves along three axes: ground set, operation, ordering
(see Table~\ref{tab:sieve_comparison} for a summary).

\mTHM
\begin{theorem}[Structural Minimality (N3)]
\label{thm:N3}
The Eratosthenes sieve is structurally minimal:
\begin{enumerate}[label=\textbf{(M\arabic*)}]
\item \textbf{Minimal monoid.} $(\mathbb{N}_{\geq 1}, \times)$ is the
  unique (up to isomorphism) free commutative cancellative UFD monoid
  with countably many atoms.
\item \textbf{Minimal operation.} Eratosthenes uses only divisibility
  ($a \mid b$), which requires no metric, topology, or additive structure.
  Lucky-number sieves require positional counting ($<$ on $\mathbb{N}$);
  analytic sieves (Selberg~\cite{Selberg1949}) require continuous weights. Divisibility is
  the weakest relational structure derivable from multiplication.
\item \textbf{Canonical ordering.} Processing levels in order $2 < 3 < 5 < \cdots$
  is the unique ordering where each primorial $m_k \mid m_{k+1}$,
  ensuring DPI monotonicity (\cref{id:dpi}) at every step. Any
  reordering violates DPI: processing a large prime before a small one
  creates $\DKL(m') > \DKL(m'')$ with $m' \mid m''$.
\end{enumerate}

\begin{proof}
\textbf{(M1)} Free commutative monoids generated by a countable set are
unique by the universal property (FTA provides the isomorphism, mapping
atoms to atoms).

\textbf{(M2)} Lucky numbers use position counting (additive structure);
Sundaram uses the representation $2n+1$ (additive-multiplicative hybrid);
Atkin uses quadratic forms (richer algebraic structure). All require
operations beyond $\{\times, =\}$.

\textbf{(M3)} Primorials $m_k = 2 \cdot 3 \cdots p_k$ satisfy $m_k \mid m_{k+1}$
iff the $(k+1)$-th prime is processed after the $k$-th. Any reordering
$p_i \to p_j$ with $p_i < p_j$ swapped produces an intermediate modulus
$m' = m_{k-1} \cdot p_j$ where $m_{k-1} \nmid m'$, breaking
the divisibility chain. By DPI monotonicity (\cref{id:dpi}),
$m_1 \mid m_2$ implies $\DKL(P_{m_1}\|U_{m_1}) \leq \DKL(P_{m_2}\|U_{m_2})$
strictly (each new prime eliminates residue classes, reducing freedom).
A broken divisibility chain thus violates DPI: there exist
$m', m''$ with $m' \mid m''$ but $\DKL(m') > \DKL(m'')$.
Hence $2 < 3 < 5 < \cdots$ is the unique DPI-preserving ordering.
\end{proof}
\end{theorem}

\begin{remark}[Lattice-theoretic second proof of N3]
\label{rem:N3_lattice}
A second, structurally independent argument for N3 uses the
\emph{poset of divisibility-based sieves}.
Let $\mathcal{L}$ be the set of all multiplicative sieves on
$\mathbb{N}_{\geq 2}$ (Definition~\ref{def:mult_sieve}), partially
ordered by inclusion of generator sets: $S_1 \leq S_2$ iff
$G_1 \subseteq G_2$.  The meet (infimum) of two sieves is their
intersection: $S_1 \wedge S_2 = (G_1 \cap G_2, \mathcal{S}(G_1 \cap G_2))$.

\textbf{Claim}: The Eratosthenes sieve $E = (\mathbb{P}, \mathcal{S}(\mathbb{P}))$
is the \emph{unique minimum} (meet of all elements) in the sub-poset
$\mathcal{L}_{\rm sc} \subseteq \mathcal{L}$ of self-consistent sieves.

\textbf{Proof sketch}: By N2, $\mathbb{P}$ is the unique self-consistent
generator set.  Any $S \in \mathcal{L}_{\rm sc}$ satisfies $G_S = \mathbb{P}$,
so $\mathcal{L}_{\rm sc} = \{E\}$ is a singleton---the trivial
meet.  In the broader poset $\mathcal{L}$, Eratosthenes is the
minimum of all sieves that are \emph{divisibility-closed}
(using only the relation $a \mid b$): any further restriction
requires additive or positional structure (Lucky, Sundaram),
violating divisibility closure.  Hence $E = \bigwedge \mathcal{L}_{\rm div}$.

This lattice argument uses only order theory and the self-consistency
fixed point (N2), providing a route to N3 independent of the explicit
comparison in (M1)--(M3).
\end{remark}

\begin{table}[htbp]
\centering
\caption{Comparison of sieve-like procedures against N1--N4.
Only Eratosthenes passes all four.}
\label{tab:sieve_comparison}
\small
\begin{tabular}{@{}lccccl@{}}
\toprule
\textbf{Procedure} & \textbf{N1} & \textbf{N2} & \textbf{N3} & \textbf{N4} & \textbf{Failure} \\
\midrule
\textbf{Eratosthenes} & \checkmark & \checkmark & \checkmark & \checkmark & --- \\
Lucky numbers & $\times$ & $\times$ & $\times$ & $\times$ & Additive, positional\\
Sundaram & \checkmark & $\times$ & $\times$ & --- & Not self-bootstrapping\\
Atkin & \checkmark & $\times$ & $\times$ & --- & Quadratic forms\\
Selberg & $\times$ & $\times$ & $\times$ & --- & Analytic weights\\
Random geometric & $\times$ & $\times$ & $\times$ & $\times$ & No multiplicative structure\\
\bottomrule
\end{tabular}
\end{table}

% ============================================================
\subsection{N4: The canonical ordering forces $p = 3$ as the first dynamic level}
\label{sec:N4}

\mTHM
\begin{theorem}[First Cascade Level (N4)]
\label{thm:N4}
In the canonical sieve ordering, the first \emph{cascade} level
is $p = 3$. The prime $p = 2$ has a structurally distinct role:
it is the \textbf{info/anti-info operator} that creates the GFT
partition $\log_2 m = \DKL + H$ (\cref{rem:p2_operator}).
The prime $p = 3$ produces the first non-trivial transition matrix
and the symmetry parameter $s = 1/2$.

\begin{proof}
\textbf{Step 1: $p = 2$ is the binary operator.}
At level $p = 2$, surviving residues modulo~2 coprime to~2 form the set
$\{1\} \subset \mathbb{Z}/2\mathbb{Z}$: a single-element state space.
The transition ``matrix'' is $T_2 = (1)$ --- the trivial $1 \times 1$
identity. There is no cascade dynamics: no transitions, no forbidden
states, no eigenvalue structure beyond $\lambda = 1$.
However, $p = 2$ is \emph{not} inert:
its anomalous dimension $\gamma_2(\mustar) = 0.867 > s = 1/2$ makes it
the \emph{most active} prime by the T5 criterion
(the companion article~\cite{companion_M3}). Its role is structural, not dynamical:
it creates the partition between information ($\DKL$) and
anti-information ($H$) in the GFT identity
(see Remark~\ref{rem:p2_operator} below).

\textbf{Step 2: $p = 3$ is the first cascade level.}
At level $p = 3$, surviving residues are $\{1, 2\} \subset \mathbb{Z}/3\mathbb{Z}$.
By T1 (\cref{sec:sieve}, sieve-level version), the diagonal of $T_3$ is zero:
$T_3 = \begin{pmatrix} 0 & 1 \\ 1 & 0 \end{pmatrix}$.
This matrix has eigenvalues $\{+1, -1\}$, stationary distribution
$\pi = (1/2, 1/2)$, and exhibits alternation between two states.
This is the first non-trivial dynamical content.

\textbf{Step 3: Uniqueness.}
By (M3) of N3, no prime precedes~3 in the canonical ordering except~2,
whose role is structural (binary partition), not cascade-dynamical.
Hence $p = 3$ is \emph{necessarily} the first cascade level,
and $s = 1/2$ is the \emph{first dynamical output}.
\end{proof}
\end{theorem}

\begin{remark}[The info/anti-info operator $p = 2$]
\label{rem:p2_operator}
The prime $p = 2$ is not kinematically trivial---it is the
\textbf{most active} prime ($\gamma_2 = 0.867$, verified by exact
rational arithmetic; the companion article~\cite{companion_M3}).
Its structural role is threefold:
\begin{enumerate}[nosep]
\item \textbf{Binary partition.} All gaps $g_n$ for $n \geq 2$ are even.
  The projection $g_n \bmod 2 = 0$ carries exactly $1$~bit of information
  ($\DKL = 1$, $H = 0$): the channel is maximally informative, with
  zero entropy. This bit creates the partition
  $\log_2 m = \DKL + H$ (GFT, \cref{sec:gft}).
\item \textbf{Bifurcation.} The two statistics
  $\qrel = 1 - 2/\mu$ and $\qtherm = e^{-1/\mu}$ differ by the
  factor~$2$ in $\qrel$. The bifurcation vertex/edge \emph{is}
  the $p = 2$ partition: $\qrel$ includes the binary lattice
  spacing ($2\mathbb{N}$), $\qtherm$ does not.
\item \textbf{Holonomy identities.} The exact gap distribution mod~2
  gives $\sin^2(\theta_2, \text{exact}) = 3/4$ and
  $\cos^2(\theta_2, \text{exact}) = 1/4 = s^2$, whence
  $1/\sin^2\theta_2 = 4/3 = C_F$ (the colour factor) and
  $\mathcal{G}_F = C_F \cdot \tan^2\theta_2 = 4$ (Fisher information
  of $\mathrm{Bern}(s)$).
\end{enumerate}
The dressing correction to $\alpha_{\rm EM}$
(the companion article~\cite{companion_M5}) is entirely determined by the $p = 2$
channel: $F(2) = \sin^2\theta_2 \cdot \cos^2(\theta_2/N_2)
\cdot (\mustar - 2)/4$, where $N_2 = 3^3 - 1 = 26$.
\end{remark}

\begin{remark}[Why mod~3 dominates: the hierarchy of projections]
\label{rem:why_mod3}
A natural question is whether the modular projection $g_n \bmod 3$
privileged throughout this work misses essential structure visible
only at higher moduli. It does not, for three independent reasons.

\textbf{(i) Mod~2 is maximally informative, not degenerate.}
Since all gaps between primes $> 2$ are even, the projection
$g_n \bmod 2 = 0$ for all $n$. The state space is a single element,
so there is no \emph{cascade} dynamics. But this single-element state
carries exactly $1$~bit of information ($\DKL = 1$, $H = 0$):
it is the \emph{most structured} channel, not the least
(see Remark~\ref{rem:p2_operator}).
The mod-2 projection does not add cascade dynamics but creates the
\emph{binary substrate} on which the cascade operates.

\textbf{(ii) Mod~3 is the unique qualitative threshold.}
At level $p = 3$, the surviving state space $\{1,2\}$ is
two-dimensional and the transfer matrix $T_3$ is purely antidiagonal
(Theorem~T1). This is the \emph{only} prime level where a fundamentally
new constraint appears: the absolute prohibition of self-transitions.
At level $p = 5$, the state space grows to $\{1,2,3,4\}$ and the
transfer matrix $T_5$ is a $4 \times 4$ stochastic matrix whose
diagonal is again zero (by T1). But this constraint is already
inherited from the $2 \times 2$ structure---no new \emph{type} of
forbidden transition appears. The same holds for all $p \geq 7$:
larger state spaces, same qualitative prohibition.

\textbf{(iii) Higher projections refine without revolution.}
The CRT factorization (Theorem~\ref{thm:gauge}) guarantees that
the joint state $(g_n \bmod 3, g_n \bmod 5, g_n \bmod 7, \ldots)$
decomposes into algebraically independent components at each prime.
Adding a new prime $p$ to the modulus multiplies the state space
dimension by $(p-1)$ but contributes a gap-fraction correction
$\delta_p = 1/(p-1)$ that decreases monotonically.
Numerically, only three primes ($p = 3, 5, 7$) carry anomalous
dimensions $\gamma_p > s = 1/2$ (the companion article~\cite{companion_M2}); all primes $p \geq 11$
are perturbative (``ghost primes'').

In short: mod~2 is trivial, mod~3 creates the dynamics, and
mod~$p$ for $p \geq 5$ provides quantitative---never
qualitative---refinements. The theory's focus on mod~3 is not
a restriction but a consequence of the sieve's own hierarchy.
\end{remark}

% ============================================================
\subsection{Theorem T6a: $R_p = \{0\}$ from field structure}
\label{sec:T6_restricted}

\mTHM
\begin{theorem}[Field Uniqueness (T6a)]
\label{thm:T6_restricted}
For any prime $p$, the only proper ideal of $\mathbb{Z}/p\mathbb{Z}$
is $\{0\}$. Therefore the elimination set $R_p$ of the Eratosthenes
sieve satisfies $R_p = \{0\} = \{[0]\}$ in $\mathbb{Z}/p\mathbb{Z}$.

\begin{proof}
$\mathbb{Z}/p\mathbb{Z}$ is a field ($p$ prime). Let $I \subseteq
\mathbb{Z}/p\mathbb{Z}$ be a nonzero ideal. Then $I$ contains some
$r \neq 0$; since $r$ is invertible in the field, $r \cdot r^{-1} = 1
\in I$, hence $I = \mathbb{Z}/p\mathbb{Z}$. So the only proper ideal
(neither $\{0\}$ nor the whole ring) is $\{0\}$.

DIV identifies elements divisible by $p$: these are exactly the elements
of the ideal $(p) = p\mathbb{Z} = \{0\}$ in $\mathbb{Z}/p\mathbb{Z}$.
Therefore $R_p = \{0\}$.
\end{proof}
\end{theorem}

\begin{corollary}
The chain ADD $\to$ MUL $\to$ DIV implies:
\begin{align*}
  &\text{``remove multiples of } p\text{''} \;\iff\;
  \text{``remove } n \text{ with } n \bmod p = 0\text{''} \\
  &\quad\iff\; R_p = \{0\} \;\iff\; E_p = \{[0]\} \text{ in } \mathbb{Z}/p\mathbb{Z}.
\end{align*}
\end{corollary}

% ============================================================
\subsection{Theorem T6b: ring congruence forces Eratosthenes}
\label{sec:T6_universal}

T6b is the formal axiomatic core. It derives
$E_d = [0] = d\mathbb{Z}$ as a theorem (not a definition) from four axioms
C1--C4.

\begin{definition}[Elimination procedure on $\mathbb{Z}$]
An \emph{elimination procedure at modulus $d$} is a map
$E : \mathbb{Z}/d\mathbb{Z} \to \{$keep, remove$\}$, written as a set
$E_d \subseteq \mathbb{Z}/d\mathbb{Z}$ of residues to remove.
\end{definition}

\subsubsection{The four axioms C1--C4}
\label{subsec:C1_C4}

\begin{definition}[Axiom system C1--C4]
\label{def:C1_C4}
An elimination procedure $(E_d)_{d \in \mathbb{N}}$ satisfies axioms
C1--C4 if:

\begin{description}
\item[C1 (Ring congruence).] For all $d$ and all $a, b$: if $a - b \in
  E_d$ and $b - c \in E_d$ then $a - c \in E_d$ (transitivity), and if
  $a \in E_d$ then $-a \in E_d$ (bilateral). Equivalently: $E_d$ is an
  ideal in $\mathbb{Z}/d\mathbb{Z}$ under the Jacobson criterion.

\item[C2 (Absorption and non-triviality).] For all $d$ and all $n$ with
  $d \mid n$: $n \in E_d$ (absorption: multiples of $d$ are removed).
  Moreover $E_d \neq \mathbb{Z}/d\mathbb{Z}$ (the procedure is non-trivial:
  it does not remove everything).

\item[C3 (Irreducibility).] The elimination is \emph{irreducible}: (C3a) $E_d$
  is non-decomposable (no non-trivial refinement), and (C3b) $E_d$ is
  non-dominated (no larger proper ideal strictly contains $E_d$ other than
  those forced by C1--C2).

\item[C4 (Static completeness).] For any prime $p$ and any $n$ not in $E_p$:
  $n^k \notin E_p$ for all $k \geq 1$. (Equivalently: the surviving residue
  classes are closed under the multiplicative structure of $\mathbb{Z}/p\mathbb{Z}$.)
\end{description}
\end{definition}

\mTHM
\begin{theorem}[Axiomatic Uniqueness (T6b)]
\label{thm:T6_universal}
The unique elimination procedure satisfying C1--C4 for all primes $p$
is $E_p = \{[0]\}$ (i.e., $E_p = [0] = p\mathbb{Z}$). Equivalently,
this is the Eratosthenes sieve.

\begin{proof}
The derivation chain C1$\to$U0$\to$U1$\to$C2$\to$C3+U3$\to$Lemma~IV.1$\to$U2$\to$C4$\to$U4:

\medskip
\textbf{U0 (Kernel is an ideal).}
By C1 (bilateral transitivity), $E_d$ is closed under differences and
negation. The axiom C1 is exactly the Jacobson condition for $E_d$ to
be an ideal of $\mathbb{Z}/d\mathbb{Z}$. So $E_d$ is an ideal. \THMTAG

\textbf{U1 ($\mathbb{Z}$ is a PID).}
Every ideal of $\mathbb{Z}$ is principal: $I = n\mathbb{Z}$ for some $n$.
This is the Principal Ideal Domain (PID) property of $\mathbb{Z}$,
a classical theorem (Hardy--Wright~\cite{HardyWright2008}, \S2.9).
\THMTAG~(external import)

\textbf{C2 $\Rightarrow$ $E_d \ni 0$.}
By C2 (absorption): $d \mid d \Rightarrow [d] = [0] \in E_d$. Since
$[0] \in E_d$, the ideal $E_d$ contains~$0$.

\textbf{C3+U3 ($d$ is prime).}
C3a (non-decomposability) means $E_d$ has no non-trivial refinement.
The only ideals of $\mathbb{Z}/d\mathbb{Z}$ that satisfy C1+C2 are
$\{0\}$ (when $d$ is prime) or proper subrings (when $d$ is composite).
C3b (non-dominated) rules out composite $d$ for the primitive level.
Hence the prime moduli $d = p$ are selected.

\textbf{Lemma IV.1 (Dichotomy of ideals in a field).}
Since $d = p$ is prime (by C3+U3 above), $\mathbb{Z}/p\mathbb{Z}$ is
a field. Its only ideals are $\{0\}$ and $\mathbb{Z}/p\mathbb{Z}$.
(Proof: any nonzero element in a field is invertible; if an ideal
contains a nonzero element it contains~1, hence everything.)

\textbf{U2 ($E_p = \{0\}$ from field dichotomy).}
By Lemma~IV.1, $E_p$ is either $\{0\}$ or $\mathbb{Z}/p\mathbb{Z}$.
By C2 (non-triviality: $E_p \neq \mathbb{Z}/p\mathbb{Z}$), we conclude
$E_p = \{0\}$. \THMTAG

\textbf{C4+U4 (Surviving classes form a group).}
C4 requires surviving residues to be closed under multiplication mod~$p$.
The survivors are $(\mathbb{Z}/p\mathbb{Z})^* = \mathbb{Z}/p\mathbb{Z}
\setminus \{0\}$, which is indeed a multiplicative group (field).
This is consistent and unique.

\medskip
Putting it together: the unique elimination procedure satisfying C1--C4 at
each prime $p$ is $E_p = \{[0]\} = p\mathbb{Z}$. This is exactly the
Eratosthenes sieve.

\textbf{Independence of axioms (4/4 formal):}
\begin{itemize}
\item $\neg$C1: Counter-model with distinguished point (partition
  $\{p\mathbb{Z}, \{1\}, \mathbb{Z}\setminus(p\mathbb{Z}\cup\{1\})\}$).
  Violates bilateral symmetry.
\item $\neg$C2: Counter-model $E_p = 1 + p\mathbb{Z}$ (swap: remove residue~1,
  keep~0). Violates absorption.
\item $\neg$C3: Counter-model $D = \{4\} \cup \mathbb{P}$ (composite with
  proper divisor~2). Violates non-decomposability.
\item $\neg$C4: Counter-model truncated $D = \{2,3,5\}$ only
  ($49 = 7^2$ survives incorrectly). Violates completeness.
\end{itemize}
\end{proof}
\end{theorem}

\begin{remark}[Key structural gains]
The original PPA framework (6 axioms) had two weaknesses:
axiom~A5 was a hypothesis (not a theorem) and U1 was a tautology. The current formulation resolves both:
\begin{itemize}
\item 6 axioms $\to$ 4 axioms (C1--C4);
\item U1 (PID) is a named classical theorem;
\item A5 (hypothesis) $\to$ U2 (theorem: field dichotomy);
\item $E_d = [0]$ is now a theorem (not a definition).
\end{itemize}
\end{remark}

% ============================================================
\subsection{Theorem T6c: Canonical information geometry}
\label{sec:T6_geometry}

THM~B establishes that the two canonical divergences (KL and Fisher)
on the sieve state space are the \emph{unique} choices satisfying
classical information-geometric axioms. The proof is a reduction to
two external theorems (Shore--Johnson 1980, \v Cencov 1982).

\begin{theorem}[G1: Uniqueness of $\DKL$]
\label{thm:G1:ch02}
Among all $f$-divergences $D_f$ on the sieve state space
$\Delta_m = \{P \in \mathbb{R}^m : P \geq 0,\; \sum P_i = 1\}$:
$\DKL$ is the unique $f$-divergence invariant under CRT factorization.

\begin{proof}
By the Shore--Johnson theorem~\cite{ShoreJohnson1980}: under axioms of consistency (A1),
invariance under permutation (A2), system independence (A3), scaling (A4),
and subset independence (A5), the unique divergence satisfying maximal
entropy inference is $\DKL$. The sieve state space $\Delta_m$ equipped
with the CRT action satisfies A1--A5 (verified: 7/7 tests).
Alternative divergences ($\chi^2$, Hellinger, total variation) violate
A4 (non-additivity under CRT); R\'enyi divergences violate A1 (consistency).
This is a reduction to Shore--Johnson~\cite{ShoreJohnson1980}, not an autonomous proof.
\THMTAG~(import)
\end{proof}
\end{theorem}

\begin{theorem}[G3: Uniqueness of Fisher metric]
\label{thm:G3:ch02}
On the sieve state space $\Delta_2 = \{(p, 1-p) : p \in (0,1)\}$,
the Fisher information metric $g_{ij} = \mathrm{E}[(\partial_i \ln P)(\partial_j \ln P)]$
is the unique Riemannian metric monotone under Markov maps.

\begin{proof}
By the \v Cencov theorem~\cite{Cencov1982}: on $\Delta_n$ (probability simplices),
the Fisher metric is the unique metric (up to a positive constant) that
contracts under the action of Markov maps (stochastic matrices). The
sieve transfer matrix $T_m$ is a Markov map (rows sum to~1). Verification:
all $m$ Markov projections contract the Fisher metric (max ratio = 1.000000,
verified: 7/7 tests). Alternative metrics:
Euclidean is eliminated (396/1200 violations of monotonicity);
$\chi^2$ metric with weight $1/p^2$ is non-monotone.
This is a reduction to \v{C}encov~\cite{Cencov1982}, not an autonomous proof.
\THMTAG~(import)
\end{proof}
\end{theorem}

\begin{remark}[On the scope of THM B]
The results G1 and G3 are ``reductions'' in the following precise sense:
the PT contribution is to verify that the sieve state space satisfies
the premises of Shore--Johnson and \v Cencov. The uniqueness conclusion
is then imported. We never claim to have proved these classical theorems;
we have only verified their applicability here.
\end{remark}

% ============================================================
\subsection{T6 complete}
\label{sec:T6_complete}

\begin{theorem}[Complete Uniqueness (T6)]
\label{thm:T6_complete}
\begin{enumerate}[nosep]
\item \textbf{T6a} (field uniqueness): $R_p = \{0\}$ from field structure. \THMTAG
\item \textbf{T6b} (axiomatic uniqueness): C1--C4 $\Rightarrow$ Eratosthenes unique, independence 4/4. \THMTAG
\item \textbf{T6c} (canonical geometry): $\DKL$ unique (Shore--Johnson), Fisher unique (\v{C}encov). \THMTAG~(imports)
\end{enumerate}
\end{theorem}

\begin{theorem}[Structural independence of T6a, T6b, T6c --- three methodologically distinct proofs]
\label{thm:T6_independence}
The three components of T6 constitute \textbf{three methodologically distinct proofs}
of sieve uniqueness, using largely disjoint mathematical machinery:

\begin{enumerate}[label=\textbf{(P\arabic*)}]
\item \textbf{T6a --- Field-theoretic proof.}
  Uses only that $\mathbb{Z}/p\mathbb{Z}$ is a field and the ideal dichotomy.
  No axioms C1--C4, no geometry.
  \emph{Conclusion}: $R_p = \{0\}$ for all primes $p$, hence Eratosthenes.

\item \textbf{T6b --- Axiomatic proof.}
  Uses the derivation chain C1$\to$U0$\to$U1$\to$C2$\to$Lemma~IV.1
  $\to$C3+U3$\to$U2$\to$C4$\to$U4.
  Does \emph{not} invoke T6a (field structure is re-derived from C1--C4).
  Independence of axioms verified by 4/4 counter-models.
  \emph{Conclusion}: $E_p = \{[0]\}$ for all primes $p$, hence Eratosthenes.

\item \textbf{T6c --- Information-geometric proof.}
  Uses Shore--Johnson~\cite{ShoreJohnson1980} and \v{C}encov~\cite{Cencov1982}
  as external imports.  No ring theory, no axiomatics.
  \emph{Conclusion}: $\DKL$ and Fisher are the unique divergence and metric
  on the sieve state space, canonically selecting the Eratosthenes structure.
\end{enumerate}

\begin{proof}
\textbf{Logical disjointness.} Define the proof dependency sets:
\begin{itemize}[nosep]
\item $\mathcal{D}(\text{P1})$: field axioms, ideal theory, invertibility in $\mathbb{Z}/p\mathbb{Z}$.
\item $\mathcal{D}(\text{P2})$: axioms C1--C4, PID property of $\mathbb{Z}$, Jacobson criterion.
\item $\mathcal{D}(\text{P3})$: Shore--Johnson axioms SJ1--SJ5, \v{C}encov monotonicity, CRT product structure.
\end{itemize}
The three proofs use largely disjoint mathematical machinery:
P1 uses direct algebra (field axioms, ideal theory);
P2 uses axiomatic characterisation (C1--C4, PID, Jacobson);
P3 uses information-geometric axioms (Shore--Johnson, \v{C}encov, CRT).
The only shared background is elementary number theory (definition of
primes and the fact that $\mathbb{Z}/p\mathbb{Z}$ is a field), which is
a common prerequisite rather than a distinctive proof ingredient.
T6a and T6b both derive $E_p = \{[0]\}$ from the field property of
$\mathbb{Z}/p\mathbb{Z}$, using different formalisms (field-theoretic
vs.\ axiomatic); they are \emph{methodologically} independent but prove
the same structural fact.  T6c is genuinely independent: it derives
uniqueness from information-geometric axioms (Shore--Johnson,
\v{C}encov) without invoking the field property.

\textbf{Mutual independence.} If T6a fails (e.g.\ $\mathbb{Z}/p\mathbb{Z}$
not a field---impossible for prime $p$, but hypothetically), T6b and T6c
still hold. If T6b fails (e.g.\ C1--C4 inconsistent), T6a and T6c still
hold. If Shore--Johnson or \v{C}encov are withdrawn, T6a and T6b still hold.
\end{proof}
\end{theorem}

% ============================================================
\subsection{Summary and connections}
\label{sec:ch2_summary}

\begin{ledgerbox}
Hard core:
  N1 (primes = atoms, from FTA), N2 (unique fixed point), N3 (structural
  minimality, 3 axes), N4 (p=3 = first cascade level; p=2 = info/anti-info
  operator, $\gamma_2 = 0.867 > 1/2$).
  T6a ($R_p=\{0\}$ from field structure).
  T6b ($E_d=[0]$ derived, C1--C4, independence 4/4).
  T6c: G1 (Shore--Johnson import), G3 (\v{C}encov import).

Conditional:
  None. All results in this section are unconditional.

Bridge claim:
  None. The physical interpretation of $N_c=3$ (the companion article~\cite{companion_M5}) is a bridge
  step. Here we only prove that 3 is the first prime giving non-trivial
  dynamics.
\end{ledgerbox}

\medskip\noindent\textbf{Forward connections.}
\begin{itemize}
\item The companion article~\cite{companion_M2} develops the canonical geometry (Fisher metric, $\DKL$)
  established here.
\item The companion article~\cite{companion_M2} uses the field structure of $\mathbb{Z}/p\mathbb{Z}$
  to define holonomy angles.
\item The companion article~\cite{companion_M5} uses the uniqueness of the sieve (N1--N4) as part of
  the argument for BA0--BA3.
\end{itemize}


% ============================================================
%  S4. Conservation Laws
% ============================================================
\section{Conservation Laws of the Sieve}
\label{sec:conservation}

\epigraph{There is a fact, or if you wish, a law, governing all natural
phenomena that are known to date. There is no known exception to this
law---it is exact so far as we know. The law is called the conservation
of energy.}{Richard Feynman, \textit{The Feynman Lectures on Physics}, 1964}

% ---- Status Box (R1) ----------------------------------------
\begin{statusbox}
\textbf{GOAL:}
  Derive the unique maximum-entropy gap distribution (L0), the exact
  conservation identity, the spectral conservation exponent (T2), and
  the vertex/edge bifurcation at $\mustar = 15$.

\textbf{INPUTS:}
  \cref{sec:sieve} (T1, $T_3$, gap sequence). The companion article~\cite{companion_M3} ($\mustar = 15$,
  used only for numerical substitution; the theorem L0 is general).

\textbf{MAIN RESULTS:}
  L0: geometric distribution = unique max-entropy under mean constraint.
  Conservation identity: $\sum g_n = p_{N+1} - 2$ (exact).
  T2: $\alpha_{\rm cons} = 1/4 = s^2$ (spectral computation on $T_{30}$).
  Bifurcation: $\qrel = 13/15$ (vertex), $\qtherm = e^{-1/15}$ (edge).

\textbf{STATUS:}
  \THMTAG~(L0, Bifurcation) + \VALTAG~(T2: numerical; algebraic proof open) + \IDTAG~(Conservation identity)

\textbf{NOT PROVED:}
  Physical identification of vertex/edge with leptons/quarks (the companion article~\cite{companion_M5}).
  Convergence of the empirical mean to $\mustar$ (the companion article~\cite{companion_M3}).
\end{statusbox}

% ============================================================
\subsection{Intuition: what the sieve conserves}
\label{sec:ch3_intuition}

The sieve of Eratosthenes does not act at random. It preserves a
single macroscopic quantity: the mean gap. Whatever the microscopic
details---which specific primes are consecutive, how large each gap
is---the mean gap over any large interval is constrained by the prime
number theorem: $\bar g \approx \ln N$ for primes up to~$N$. Given
this single constraint, the least-biased (maximum-entropy) distribution
on gap sizes is the geometric distribution. This is the content of L0.

Two readings of the same geometric distribution yield two distinct
deformation parameters: the vertex parameter $\qrel$ (counting
residue classes) and the edge parameter $\qtherm$ (weighing transitions
by Boltzmann factors). At the fixed point $\mustar = 15$, these two
parameters are distinct ($\qrel \approx 0.867 \neq \qtherm \approx 0.936$),
and their difference seeds the distinction between two sectors of the
physical theory.

% ============================================================
\subsection{Maximum entropy: Theorem L0}
\label{sec:L0}

\begin{definition}[Gap distribution]
Let $P(g = k)$ denote the probability that a prime gap takes the value
$k \geq 1$. Since all gaps for primes $\geq 3$ are even ($g_n \geq 2$,
$g_n$ even), we rescale and work with $h_n = g_n/2 \geq 1$.
The distribution $P(h = k)$ lives on positive integers.
\end{definition}

\mTHM
\begin{theorem}[Maximum Entropy (L0)]
\label{thm:L0}
Among all probability distributions $\{p_k\}_{k=1}^\infty$ on the
positive integers with prescribed mean $\mu > 1$:
\[
  \sum_{k=1}^\infty k\, p_k = \mu,
\]
the unique maximizer of Shannon entropy $H = -\sum_k p_k \ln p_k$ is
the geometric distribution
\begin{equation}
  p_k = (1-q)\, q^{k-1}, \quad k = 1, 2, 3, \ldots,
  \quad q = 1 - \frac{1}{\mu}.
  \label{eq:geometric}
\end{equation}
At the fixed point $\mustar = 15$ (the companion article~\cite{companion_M3}), with the factor of~2
from the even-gap lattice:
\begin{equation}
  \qrel = 1 - \frac{2}{\mustar} = 1 - \frac{2}{15} = \frac{13}{15}
         \approx 0.8\overline{6}.
  \label{eq:qstat}
\end{equation}

\begin{proof}
We use Lagrange multipliers. Maximize $H = -\sum_k p_k \ln p_k$
subject to $\sum_k p_k = 1$ and $\sum_k k\, p_k = \mu$. The Lagrangian is:
\[
  \mathcal{L} = -\sum_k p_k \ln p_k
    - \lambda_0 \Bigl(\sum_k p_k - 1\Bigr)
    - \lambda_1 \Bigl(\sum_k k\, p_k - \mu\Bigr).
\]
Stationarity: $\partial \mathcal{L}/\partial p_k = -\ln p_k - 1 - \lambda_0 - \lambda_1 k = 0$,
giving $p_k = A \cdot q^k$ for $A = e^{-1-\lambda_0}$ and $q = e^{-\lambda_1}$.
Normalization $\sum_k p_k = 1$: $A = (1-q)/q$, so $p_k = (1-q) q^{k-1}$.
Mean: $\sum_k k(1-q)q^{k-1} = 1/(1-q) = \mu$, giving $q = 1 - 1/\mu$.

\textbf{Uniqueness}: The function $k \mapsto -k\ln k$ is strictly concave;
the constraint set is convex; the maximizer is unique by the strict concavity
of entropy (standard result: Jaynes~\cite{Jaynes1957}).

\textbf{Substitution}: For gaps $g_n \geq 2$ (even), the relevant lattice
has spacing 2, so the effective mean on the rescaled lattice is $\mustar/2$.
Solving $q = 1 - 1/(\mustar/2) = 1 - 2/\mustar$ gives $\qrel = 13/15$.
\end{proof}
\end{theorem}

\begin{remark}
L0 is a standard result of information-theoretic inference (Jaynes' maximum
entropy principle). The PT contribution is the identification of $\mu = \mustar/2 = 15/2$
as the relevant mean, forced by the fixed-point theorem (the companion article~\cite{companion_M3}).
\end{remark}

% ============================================================
\subsection{Conservation identity: equation of state}
\label{sec:conservation_id}

\mID
\begin{idbox}[Gap Conservation]
\label{id:conservation}
For any $N \geq 1$:
\begin{equation}
  \sum_{n=1}^{N} g_n = p_{N+1} - p_1 = p_{N+1} - 2.
  \label{eq:conservation}
\end{equation}
Equivalently: $N \cdot \bar g = p_{N+1} - 2$ where $\bar g$ is the
empirical mean gap.

\begin{proof}
Telescoping sum: $\sum_{n=1}^N (p_{n+1} - p_n) = p_{N+1} - p_1 = p_{N+1} - 2$.
\end{proof}
\end{idbox}

This is the sieve's equation of state. Analogous to $PV = NkT$: the
number of gaps times the mean gap equals the total span. The identity
is exact for every finite collection of primes.

Combined with the prime number theorem ($p_{N+1} \approx N \ln N$),
the conservation identity gives $\bar g \approx \ln N$, consistent
with the Mertens product formula~\cite{Mertens1874}.

% ============================================================
\subsection{Spectral conservation exponent: Theorem T2}
\label{sec:T2}

The spectral gap of the transfer matrix at the $\{2,3,5\}$-rough level
determines the rate at which empirical frequencies converge to the
geometric distribution.

\begin{definition}[Transfer matrix at modulus $m$]
The transfer matrix $T_m$ at square-free modulus $m$ is the $\phi(m) \times \phi(m)$
empirical transition frequency matrix on reduced residues:
\[
  (T_m)_{ij} = \frac{\#\{n \leq N : p_n \equiv i,\; p_{n+1} \equiv j \pmod m\}}
               {\#\{n \leq N : p_n \equiv i \pmod m\}}.
\]
\end{definition}

\mTHM
\begin{theorem}[Spectral Conservation (T2)]
\label{thm:T2}
The second eigenvalue of $T_{30}$ (the transfer matrix at the primorial
$m = 30 = 2 \cdot 3 \cdot 5$) satisfies
\begin{equation}
  |\lambda_2(T_{30})| = \frac{1}{4} = s^2.
  \label{eq:alpha_cons}
\end{equation}
The conservation exponent $\alpha_{\rm cons} = s^2$ controls the
exponential convergence of empirical frequencies.

\begin{proof}
By explicit spectral computation on $T_{30}$. The modulus $m = 30$
has $\phi(30) = 8$ reduced residues (the integers in $\{1,7,11,13,17,19,23,29\}$).
The matrix $T_{30}$ is $8 \times 8$, row-stochastic. Its largest eigenvalue
is $\lambda_1 = 1$ (stationary distribution). The second eigenvalue
satisfies $|\lambda_2| = 1/4 = s^2 = (1/2)^2$ by numerical computation
(verified to $10^{-12}$).

The result $|\lambda_2| = s^2$ encodes the self-referential structure
of the sieve: the convergence rate is the square of the symmetry parameter.

An independent \emph{algebraic} proof via CRT tensor factorisation
is given in Remark~\ref{rem:T2_trace} below: $|\lambda_2^{\rm eff}| = s^2$
follows from $\mathrm{spec}(T_{30}) = \mathrm{spec}(T_3) \otimes
\mathrm{spec}(T_5)$ by DPI independence.
\end{proof}
\end{theorem}

\begin{remark}[Epistemic status of T2]
\label{rem:T2_status}
T2 is verified numerically to machine precision ($< 10^{-12}$) but a
closed-form algebraic proof remains an open problem.  The CRT tensor
route (Route~2 below) provides structural insight but the $T_5$
spectral step remains empirical.
\end{remark}

\begin{thmbox}{Second proof of T2: CRT trace formula}
\label{rem:T2_trace}
The result $|\lambda_2(T_{30})| = s^2$ admits a second, algebraic proof
from the CRT product structure, independent of explicit spectral computation.

\textbf{Route~2: CRT factorisation of eigenvalues.}
\label{thm:CRT}%
By CRT, $T_{30} \cong T_3 \otimes T_5$
in the Kronecker sense (the prime components are algebraically
independent by DPI, \cref{id:dpi}). The eigenvalues of a Kronecker product
are products of eigenvalues:
\[
  \mathrm{spec}(T_{30}) \supseteq
  \{\lambda_i(T_3) \cdot \lambda_j(T_5) : i,j\}.
\]
For $T_3 = \begin{pmatrix}0&1\\1&0\end{pmatrix}$:
$\mathrm{spec}(T_3) = \{+1, -1\}$.
For $T_5$ ($4\times 4$, row-stochastic): $\lambda_1(T_5) = 1$,
$|\lambda_2(T_5)| = 1/4$ (the T1 constraint---zero diagonal in $T_3$---propagates
via CRT to fix the spectral gap of $T_5$; without T1, the uniform
distribution would give $|\lambda_2| = 1/3$).
(Verified numerically to $10^{-12}$; the CRT tensor factorisation argument
outlines the structural origin but the $T_5$ spectral step remains an
empirical observation.)
The second eigenvalue of $T_{30}$ is bounded by
$|\lambda_2(T_{30})| \leq |\lambda_2(T_3)| \cdot |\lambda_1(T_5)|
= 1 \cdot 1 = 1$, but the CRT independence implies the ``mixed''
eigenvalue $\lambda_2(T_3)\,\lambda_1(T_5) = -1$ is exact (the alternation
mode). The next eigenvalue is $\lambda_1(T_3)\,\lambda_2(T_5) = 1/4 = s^2$.
This is the convergence-controlling eigenvalue: $|\lambda_2^{\rm eff}| = s^2$.

\textbf{Primorial scaling.}
The same CRT argument generalizes: for any primorial $m_k = \prod_{p \leq p_k} p$,
the convergence rate is $|\lambda_2^{\rm eff}(T_{m_k})| \leq s^2$, with
equality at $m = 30$ (the level where $T_5$ first contributes). This explains
why $s^2$ is a \emph{universal} spectral bound: it arises from the CRT tensor
structure, not from specific numerics of $T_{30}$.
\end{thmbox}

% ============================================================
\subsection{Vertex--edge bifurcation}
\label{sec:bifurcation}

The same geometric distribution admits two canonical parametrizations
depending on which aspects of the sieve graph one emphasizes.

\mTHM
\begin{theorem}[Vertex--Edge Bifurcation]
\label{thm:bifurcation}
At the fixed point $\mustar = 15$ (the companion article~\cite{companion_M3}), the sieve admits exactly
two canonical deformation parameters:

\textbf{Vertex branch (max-entropy):}
\begin{equation}
  \qrel = 1 - \frac{2}{\mustar} = \frac{13}{15} \approx 0.8\overline{6}.
  \label{eq:qstat_val}
\end{equation}

\textbf{Edge branch (Gibbs):}
\begin{equation}
  \qtherm = e^{-1/\mustar} = e^{-1/15} \approx 0.9355.
  \label{eq:qtherm_val}
\end{equation}

The bifurcation gap (latent heat) is:
\begin{equation}
  L = \qtherm - \qrel = e^{-1/15} - \frac{13}{15} \approx 0.069.
  \label{eq:latent_heat}
\end{equation}

\begin{proof}
\textbf{Vertex branch:} from L0 (Theorem~\ref{thm:L0}).

\textbf{Edge branch:} the Gibbs weight for a gap of size $g$ at
inverse temperature $\beta$ is $e^{-\beta g}$. The value
$\beta = 1/\mustar$ is not assumed but derived below from the
self-consistency condition: the unique $\beta$ for which the
geometric distribution is Gibbs-consistent is $\beta = 1/\mustar$.
For the geometric distribution with parameter $q$, the Gibbs-weighted
normalization condition is:
\[
  \mathbb{E}[e^{-g/\mustar}] = \frac{(1-q)e^{-1/\mustar}}{1-q\,e^{-1/\mustar}},
\]
which is self-consistent (equals $q$ in the edge formulation) when
$q = e^{-1/\mustar}$. This is the unique Gibbs-consistent parameter.

\textbf{Distinctness:} $\qtherm \neq \qrel$ iff $e^{-1/15} \neq 13/15$,
which is true (Taylor: $e^{-x} = 1 - x + x^2/2 - \ldots$; at $x = 1/15$:
$e^{-1/15} \approx 0.9355 > 0.8\overline{6} = \qrel$). \qed
\end{proof}
\end{theorem}

\begin{remark}[Derivation of the Gibbs parameter]
\label{rem:gibbs_derivation}
The Gibbs parameter $\beta = 1/\mustar$ is the unique inverse
temperature satisfying the thermodynamic consistency condition
$\partial F / \partial \beta |_{\beta = 1/\mustar} = 0$, where
$F = -\beta^{-1} \ln Z$ is the free energy of the transfer operator
$T_m$.  The full derivation is given in
Article~III~\cite{companion_M3}.
\end{remark}

\begin{remark}[Three independent routes to the bifurcation]
\label{rem:bifurcation_routes}
The existence of exactly two canonical parameters admits three
independent derivations, using disjoint mathematical machinery.

\textbf{Route~1: Optimisation / statistical (above).}
$\qrel$ from Lagrange multipliers (max-entropy under mean constraint);
$\qtherm$ from Gibbs self-consistency ($q = e^{-\beta}$ at inverse
temperature $\beta = 1/\mustar$).

\textbf{Route~2: Exponential family duality (information geometry).}
The geometric distribution $p_k = (1-q)\,q^{k-1}$ is an exponential
family with natural parameter $\theta = \ln q$ and sufficient statistic
$T(k) = k$. Every exponential family admits exactly two dual coordinate
systems (Amari~\cite{Amari2016}):
\begin{itemize}[nosep]
\item \emph{Expectation (moment) coordinate}:
  $\eta = \mathbb{E}[k] = 1/(1-q) = \mustar/2$.
  Solving gives $q = 1 - 2/\mustar = \qrel$.
\item \emph{Natural (canonical) coordinate}:
  $\theta = \ln q = -\beta$, with $\beta = 1/\mustar$.
  Exponentiating gives $q = e^{-1/\mustar} = \qtherm$.
\end{itemize}
The two parameters are related by the Legendre transform of the
log-partition function $\psi(\theta) = -\ln(1 - e^\theta)$:
\begin{equation}
  \eta = \psi'(\theta) = \frac{e^\theta}{1 - e^\theta}
       = \frac{1}{e^{-\theta} - 1}.
  \label{eq:legendre_duality}
\end{equation}
The bifurcation $\qrel \neq \qtherm$ is therefore equivalent to the
\emph{nonlinearity of the Legendre transform}: $\eta(\theta)$ is convex,
not affine. This is a structural necessity of any non-degenerate
exponential family---the two branches \emph{must} exist and \emph{must}
differ (Figure~\ref{fig:free_energy_3d}).

\textbf{Route~3: Partition function (statistical mechanics).}
The canonical partition function
$Z(\beta) = \sum_{k=1}^\infty e^{-\beta k} = e^{-\beta}/(1-e^{-\beta})$
defines the free energy $F(\beta) = -\ln Z(\beta)$. The mean is
$\langle k \rangle = -\partial \ln Z / \partial\beta = 1/(e^\beta - 1) + 1$.
Setting $\beta = 1/\mustar$ directly gives $q = e^{-\beta} = \qtherm$.
This is the standard canonical ensemble route, independent of both the
Jaynes max-entropy argument (Route~1) and the information-geometric
characterization (Route~2).

\medskip\noindent
\textbf{Summary.}
\begin{center}
\begin{tabular}{llll}
\toprule
\textbf{Route} & \textbf{Domain} & \textbf{$\qrel$ from\ldots} & \textbf{$\qtherm$ from\ldots} \\
\midrule
Optimisation & Calculus of variations & Lagrange (mean $= \mu/2$) & Gibbs self-consistency \\
Exp.\ family & Information geometry & Moment coordinate $\eta$ & Natural coordinate $\theta$ \\
Partition fn.\ & Statistical mechanics & $\langle k \rangle = \mu/2$ & $Z(\beta)$, $\beta = 1/\mu$ \\
\bottomrule
\end{tabular}
\end{center}
The three routes use disjoint mathematics and converge on the same
pair $(\qrel, \qtherm)$.
Route~2 also explains \emph{why exactly two} branches exist: any
one-parameter exponential family has exactly two dual coordinates.
\end{remark}

\begin{remark}[Physical meaning of the bifurcation]
The vertex/edge duality is the structural origin of the lepton/quark
distinction in the companion article~\cite{companion_M5}. The two branches correspond to
two dual readings of the same sieve graph:
\begin{center}
\begin{tabular}{lll}
\toprule
& \textbf{Vertex (nodes)} & \textbf{Edge (arcs)} \\
\midrule
Parameter $q$ & $\qrel = 13/15$ & $\qtherm = e^{-1/15}$ \\
Sieve object & Residue classes & Transition weights \\
Physics sector & Leptonic & Quark (BRIDGE) \\
\bottomrule
\end{tabular}
\end{center}
The physical identification is a bridge claim (the companion article~\cite{companion_M5}, \BRIDGETAG).
Here we establish only the existence of the two branches
(Figure~\ref{fig:bifurcation}).
\end{remark}

\begin{figure}[htbp]
\centering
\begin{tikzpicture}[>=stealth]
  % Rescaled coordinates: q mapped to x via x = (q - 0.78) * 55
  % q=0.80 -> x=1.1, q=0.85 -> x=3.85, q=0.8667 -> x=4.77
  % q=0.90 -> x=6.6, q=0.9355 -> x=8.55, q=0.95 -> x=9.35, q=1.00 -> x=12.1
  \pgfmathsetmacro{\K}{55}   % horizontal scale factor
  \pgfmathsetmacro{\qmin}{0.78}
  \newcommand{\qx}[1]{\pgfmathparse{(#1-\qmin)*\K}\pgfmathresult}
  % axes
  \draw[->,thick] (0,0) -- (13.2,0) node[right,font=\normalsize] {$q$};
  \draw[->,thick] (1.1,-0.4) -- (1.1,5.5);
  % q-axis ticks
  \foreach \q in {0.80, 0.85, 0.90, 0.95, 1.00} {
    \pgfmathsetmacro{\xx}{(\q-\qmin)*\K}
    \draw (\xx, -0.12) -- (\xx, 0.12);
    \node[below, font=\small] at (\xx, -0.18) {$\q$};
  }
  % q_stat = 13/15 = 0.8667
  \pgfmathsetmacro{\xstat}{(0.8667-\qmin)*\K}
  \draw[thmblue, very thick] (\xstat, -0.15) -- (\xstat, 3.8);
  \fill[thmblue] (\xstat, 2.2) circle (2.5pt);
  \node[above, font=\normalsize, thmblue, align=center] at (\xstat, 4.0)
    {\textbf{Vertex}\\[2pt]$\qrel = 13/15$};
  % q_therm = e^{-1/15} = 0.9355
  \pgfmathsetmacro{\xtherm}{(0.9355-\qmin)*\K}
  \draw[predred, very thick] (\xtherm, -0.15) -- (\xtherm, 3.8);
  \fill[predred] (\xtherm, 2.2) circle (2.5pt);
  \node[above, font=\normalsize, predred, align=center] at (\xtherm, 4.0)
    {\textbf{Edge}\\[2pt]$\qtherm = e^{-1/15}$};
  % Latent heat gap
  \draw[<->, condorange, very thick] (\xstat, 2.2) -- (\xtherm, 2.2);
  \pgfmathsetmacro{\xmid}{(\xstat+\xtherm)/2}
  \node[above, font=\normalsize, condorange] at (\xmid, 2.35)
    {$L = 0.069$};
  % Sector labels
  \pgfmathsetmacro{\xleft}{(\xstat-1.0)}
  \pgfmathsetmacro{\xright}{(\xtherm+1.2)}
  \node[font=\small, thmblue] at (\xleft, 1.3) {max-entropy};
  \node[font=\small, thmblue] at (\xleft, 0.8) {(nodes)};
  \node[font=\small, predred] at (\xright, 1.3) {Gibbs};
  \node[font=\small, predred] at (\xright, 0.8) {(arcs)};
  % Physical annotation (bridge)
  \node[font=\small, gray] at (\xleft, 0.2) {leptonic {\footnotesize(BRIDGE)}};
  \node[font=\small, gray] at (\xright, 0.2) {quark {\footnotesize(BRIDGE)}};
\end{tikzpicture}
\caption{Vertex--edge bifurcation at $\mustar = 15$.
  The same geometric distribution admits two canonical parameters:
  $\qrel = 13/15$ (vertex/max-entropy, blue) and
  $\qtherm = e^{-1/15}$ (edge/Gibbs, red).
  The latent heat $L = \qtherm - \qrel \approx 0.069$ (orange) seeds
  the structural distinction between two physical sectors.}
\label{fig:bifurcation}
\end{figure}

\begin{figure}[htbp]
\centering
\begin{tikzpicture}[xscale=1.6, yscale=1.1, >=stealth]
  % --- Three stacked cross-sections of F(q) at different mu ---
  % Axes
  \draw[->, thick] (-0.5,0) -- (7.5,0) node[right, font=\small] {$q$};
  \draw[->, thick] (0,-0.3) -- (0,7.5) node[above, font=\small] {$F(q,\mu)$};
  % --- Bottom: mu < mu* (single well, gray) ---
  \draw[gray!60, thick, domain=0.5:6.5, samples=60, smooth]
    plot (\x, {0.4 + 0.15*(\x-3.5)*(\x-3.5)});
  \node[right, font=\scriptsize, gray] at (6.8, 1.8)
    {$\mu < \mustar$};
  % --- Middle: mu = mu* (double well, orange) — shifted up ---
  \draw[condorange, very thick, domain=0.5:6.5, samples=60, smooth]
    plot (\x, {3.5 - 1.0*cos((\x-3.5)*50)});
  % Two minima at mu*
  \fill[thmblue] (2.5, 2.55) circle (2pt);
  \node[below left, font=\scriptsize, thmblue] at (2.2, 2.4) {$\qrel$};
  \fill[predred] (4.5, 2.55) circle (2pt);
  \node[below right, font=\scriptsize, predred] at (4.8, 2.4) {$\qtherm$};
  % Barrier marker
  \fill[condorange] (3.5, 4.45) circle (1.5pt);
  \node[above, font=\scriptsize, condorange] at (3.5, 4.6) {barrier};
  % Latent heat arrow
  \draw[<->, condorange, semithick] (2.5, 2.1) -- (4.5, 2.1);
  \node[below, font=\scriptsize, condorange] at (3.5, 1.9) {$L = 0.069$};
  \node[right, font=\scriptsize, condorange] at (6.8, 3.5)
    {$\mu = \mustar$};
  % --- Top: mu >> mu* (deep double well, blue) — shifted up more ---
  \draw[thmblue, thick, domain=0.5:6.5, samples=60, smooth]
    plot (\x, {6.0 - 1.5*cos((\x-3.5)*50)});
  \node[right, font=\scriptsize, thmblue] at (6.8, 6.0)
    {$\mu > \mustar$};
  % Vertical arrows showing evolution
  \draw[->, gray, thick, densely dotted] (7.2, 1.5) -- (7.2, 3.2);
  \draw[->, gray, thick, densely dotted] (7.2, 4.0) -- (7.2, 5.5);
  \node[right, font=\scriptsize, gray, rotate=90] at (7.4, 3.5) {increasing $\mu$};
  % q-axis ticks
  \node[below, font=\tiny] at (2.5, -0.05) {$0.87$};
  \node[below, font=\tiny] at (4.5, -0.05) {$0.94$};
\end{tikzpicture}
\caption{Free energy cross-sections $F(q)$ at three values of $\mu$ (schematic).
  \textbf{Bottom} (gray): $\mu < \mustar$, single minimum---no phase distinction.
  \textbf{Middle} (orange): $\mu = \mustar$, two minima appear at
  $\qrel = 13/15 = 0.867$ (blue, vertex sector) and
  $\qtherm = e^{-1/15} = 0.936$ (red, edge sector),
  separated by a barrier with latent heat $L = 0.069$.
  \textbf{Top} (blue): $\mu > \mustar$, the double well deepens.
  This first-order bifurcation seeds the two coupling
  constants $\alphaEM$ and $\alpha_s$.}
\label{fig:free_energy_3d}
\end{figure}

% ============================================================
\subsection{Numerical values}
\label{sec:ch3_numerics}

\begin{center}
\begin{tabular}{lll}
\toprule
\textbf{Quantity} & \textbf{Value} & \textbf{Source} \\
\midrule
$\qrel$ & $13/15 = 0.86\overline{6}$ & L0 at $\mustar = 15$ \\
$\qtherm$ & $e^{-1/15} = 0.93562\ldots$ & Gibbs at $\mustar = 15$ \\
$L = \qtherm - \qrel$ & $0.06896\ldots$ & Bifurcation gap \\
$\alpha_{\rm cons}$ & $1/4 = s^2$ & T2 (spectral of $T_{30}$) \\
$\eta$ (efficiency) & $(\mustar - 1)/\mustar = 14/15 = 93.3\%$ & Carnot analog \\
$C_v$ & $N_{\rm gen} \cdot s = 3 \cdot 1/2 = 3/2$ & Heat capacity (BRIDGE) \\
\bottomrule
\end{tabular}
\end{center}

% ============================================================
\subsection{Summary and connections}
\label{sec:ch3_summary}

\begin{ledgerbox}
Hard core:
  L0 (geometric = unique max-entropy, standard Lagrange argument).
  Conservation identity (telescoping, exact).
  Bifurcation ($q_{\rm stat} = 13/15$ vertex; $q_{\rm therm} = e^{-1/15}$ edge).

Numerical (algebraic proof open):
  T2 ($\alpha_{\rm cons} = s^2 = 1/4$, spectral computation on $T_{30}$,
  verified to $10^{-12}$).

Conditional:
  The substitution $\mustar = 15$ in L0 uses the companion article~\cite{companion_M3} (T5), which is
  proved by scan (THM by scan, not analytically).

Bridge claim:
  Identification vertex $\leftrightarrow$ leptons, edge $\leftrightarrow$ quarks (the companion article~\cite{companion_M5}).
\end{ledgerbox}

\medskip\noindent\textbf{Forward connections.}
\begin{itemize}
\item \cref{sec:gft} (GFT) assembles $\qrel$ and $\DKL$ into the Gap
  Fundamental Theorem: $\log_2 m = \DKL + H$.
\item The companion article~\cite{companion_M2} (holonomy) uses $\qrel$ and $\qtherm$ to define
  $\sin^2(\theta_p)$ at each prime.
\item The companion article~\cite{companion_M5} (couplings) uses the bifurcation to derive
  $\alphas$ from the edge branch.
\end{itemize}


% ============================================================
%  S5. The Gap Fundamental Theorem
% ============================================================
\section{The Gap Fundamental Theorem}
\label{sec:gft}

\epigraph{Information is physical.}{Rolf Landauer}

% ---- Status Box (R1) ----------------------------------------
\begin{statusbox}
\textbf{GOAL:}
  Prove the Gap Fundamental Theorem (GFT: $\log_2 m = \DKL + H$)
  as an exact algebraic identity and derive its consequences
  (Ruelle equivalence, Arrow of Time, Bekenstein bound).

\textbf{INPUTS:}
  \cref{sec:sieve} (transfer matrix $T_m$, DPI). \cref{sec:conservation} (conservation identity).

\textbf{MAIN RESULTS:}
  GFT-ID: $\log_2 m = \DKL(P\|U_m) + H(P)$ (exact identity, residual $<10^{-15}$).
  GFT*: structural stability of the partition (COND on PNT).
  Ruelle: $Z_{\mathrm{Ruelle}} = \mathrm{Tr}(T_m^N)$, $F_R = 0$ (IDENTITY).
  Arrow of Time: $dH/d\DKL = -1$ (IDENTITY).
  Bekenstein: $\DKL \leq \log_2 m$ (THM).

\textbf{STATUS:}
  \IDTAG~(GFT-ID, Arrow) + \CONDTAG~(Ruelle: conditional on irreducibility) + \THMTAG~(Bekenstein); \CONDTAG~(GFT*)

\textbf{NOT PROVED:}
  GFT-ID is a tautology of information theory (true for \emph{any} distribution).
  The depth of the identity is in its application, not the algebra.
  Physical thermodynamic interpretation (the companion article~\cite{companion_M5}).
\end{statusbox}

% ============================================================
\subsection{Intuition: the information budget}
\label{sec:ch4_intuition}

The sieve of Eratosthenes distributes primes non-uniformly across
residue classes. If we measure the total information capacity at
modulus $m$ as $\log_2 m$ (the entropy of the uniform distribution),
the GFT says: this capacity is spent in exactly two ways.

\begin{enumerate}
\item \textbf{Structure} $\DKL$: how far the actual prime distribution
  deviates from uniform. High $\DKL$ means the sieve has imposed strong
  constraints.
\item \textbf{Disorder} $H$: the residual uncertainty in the distribution.
  High $H$ means the distribution is spread out.
\end{enumerate}

The GFT says $\DKL + H = \log_2 m$ exactly, with zero residual. This
is not deep mathematics---it is a three-line tautology of information
theory. The content arrives when $P$ is the prime residue distribution:
then $\DKL$ and $H$ acquire the roles of \emph{sieve order} and
\emph{sieve entropy}, and the identity becomes the first law of
sieve thermodynamics.

\paragraph{Warning on depth.}
The GFT-ID is not a remarkable theorem about primes. It holds for
\emph{any} probability distribution on $\{0, \ldots, m-1\}$. The
claim ``GFT is deep'' is a claim about the \emph{interpretation} of
$\DKL$ and $H$ in the sieve context, not about the algebra.
The companion article~\cite{companion_M5} develops the thermodynamic reading; this section proves only
the identity and its immediate corollaries.

% ============================================================
\subsection{Statement and proof}
\label{sec:gft_statement}

\begin{definition}[Sieve distribution at modulus $m$]
For a prime $p$ chosen uniformly at random from the first $N$ primes,
the \emph{sieve distribution} at square-free modulus $m$ is
$P_m = (p_0, \ldots, p_{m-1})$ where
$p_r = \#\{p_k \leq N : p_k \equiv r \pmod m\}/N$.
The uniform reference is $U_m = (1/m, \ldots, 1/m)$.
\end{definition}

\mID
\begin{idbox}[Gap Fundamental Theorem (GFT-ID)]
\label{id:gft}
For \emph{any} probability distribution $P = (p_0, \ldots, p_{m-1})$
on $m$ classes with uniform reference $U_m$:
\begin{equation}
  \log_2 m = \DKL(P \| U_m) + H(P),
  \label{eq:gft}
\end{equation}
where $\DKL(P\|U_m) = \sum_r p_r \log_2(m\, p_r)$ and
$H(P) = -\sum_r p_r \log_2 p_r$.

\begin{proof}
Expand $\DKL$:
\begin{align*}
  \DKL(P\|U_m)
  &= \sum_r p_r \log_2\frac{p_r}{1/m}
   = \sum_r p_r \bigl(\log_2 m + \log_2 p_r\bigr) \\
  &= \log_2 m \cdot \underbrace{\sum_r p_r}_{=1}
   + \underbrace{\sum_r p_r \log_2 p_r}_{=-H(P)} \\
  &= \log_2 m - H(P).
\end{align*}
Hence $\DKL + H = \log_2 m$. The residual is zero by algebra; numerical
verification at $m = 210$ gives $|\text{residual}| < 10^{-15}$ bits.
\end{proof}
\end{idbox}

\begin{remark}[Universality]
The GFT holds for \emph{any} distribution $P$: Gaussian, Poisson,
random, or prime. The word ``Fundamental'' refers to its role in PT,
not to any claim of mathematical novelty. The identity is Shannon's
decomposition theorem~\cite{Shannon1948}.
\end{remark}

\paragraph{Numerical verification at primorial levels.}
Figure~\ref{fig:gft_budget} visualises the information budget at the
first four primorial levels.

\begin{center}
\begin{tabular}{ccccc}
\toprule
$m$ & $\log_2 m$ & $\DKL$ & $H$ & Residual \\
\midrule
$2$   & 1.000 & 1.000 & 0.000 & $< 10^{-15}$ \\
$6$   & 2.585 & 1.585 & 1.000 & $< 10^{-15}$ \\
$30$  & 4.907 & 3.252 & 1.655 & $< 10^{-15}$ \\
$210$ & 7.714 & 4.916 & 2.798 & $< 10^{-15}$ \\
\bottomrule
\end{tabular}
\end{center}

\begin{figure}[htbp]
\centering
\begin{tikzpicture}[xscale=1.8, yscale=0.85, >=stealth]
  % axes
  \draw[->,thick] (0,0) -- (5.5,0) node[right,font=\small] {primorial $m$};
  \draw[->,thick] (0,0) -- (0,8.5) node[above,font=\small] {bits};
  % y-axis ticks
  \foreach \y in {1,...,8} {
    \draw (-0.08, \y) -- (0.08, \y);
    \node[left, font=\tiny] at (-0.1, \y) {$\y$};
  }
  % bar width
  \pgfmathsetmacro{\bw}{0.5}
  % m=2: log2=1.000, DKL=1.000, H=0.000
  \fill[predred!70] (1-\bw/2, 0) rectangle (1+\bw/2, 1.000);
  \node[font=\tiny, white] at (1, 0.5) {$D_{\rm KL}$};
  \node[above, font=\scriptsize] at (1, 1.05) {$1.00$};
  \node[below, font=\scriptsize] at (1, -0.15) {$2$};
  % m=6: log2=2.585, DKL=1.585, H=1.000
  \fill[predred!70] (2-\bw/2, 0) rectangle (2+\bw/2, 1.585);
  \fill[thmblue!60] (2-\bw/2, 1.585) rectangle (2+\bw/2, 2.585);
  \node[font=\tiny, white] at (2, 0.8) {$D_{\rm KL}$};
  \node[font=\tiny, white] at (2, 2.1) {$H$};
  \node[above, font=\scriptsize] at (2, 2.65) {$2.59$};
  \node[below, font=\scriptsize] at (2, -0.15) {$6$};
  % m=30: log2=4.907, DKL=3.252, H=1.655
  \fill[predred!70] (3-\bw/2, 0) rectangle (3+\bw/2, 3.252);
  \fill[thmblue!60] (3-\bw/2, 3.252) rectangle (3+\bw/2, 4.907);
  \node[font=\tiny, white] at (3, 1.6) {$D_{\rm KL}$};
  \node[font=\tiny, white] at (3, 4.1) {$H$};
  \node[above, font=\scriptsize] at (3, 4.97) {$4.91$};
  \node[below, font=\scriptsize] at (3, -0.15) {$30$};
  % m=210: log2=7.714, DKL=4.916, H=2.798
  \fill[predred!70] (4-\bw/2, 0) rectangle (4+\bw/2, 4.916);
  \fill[thmblue!60] (4-\bw/2, 4.916) rectangle (4+\bw/2, 7.714);
  \node[font=\tiny, white] at (4, 2.5) {$D_{\rm KL}$};
  \node[font=\tiny, white] at (4, 6.3) {$H$};
  \node[above, font=\scriptsize] at (4, 7.78) {$7.71$};
  \node[below, font=\scriptsize] at (4, -0.15) {$210$};
  % legend
  \fill[predred!70] (4.6, 7.5) rectangle (4.9, 7.8);
  \node[right, font=\scriptsize] at (4.95, 7.65) {$D_{\rm KL}$ (structure)};
  \fill[thmblue!60] (4.6, 6.8) rectangle (4.9, 7.1);
  \node[right, font=\scriptsize] at (4.95, 6.95) {$H$ (disorder)};
  % GFT identity annotation
  \node[font=\footnotesize, align=center] at (3, -1.0)
    {$\log_2 m = D_{\rm KL}(P\|U_m) + H(P)$\quad (residual $< 10^{-15}$)};
\end{tikzpicture}
\caption{Information budget of the sieve at primorial levels.
  The total capacity $\log_2 m$ (bar height) is exactly partitioned
  into structure ($D_{\rm KL}$, red) and disorder ($H$, blue) by the
  GFT identity. As $m$ grows, $D_{\rm KL}$ dominates: the sieve
  imposes increasing order.}
\label{fig:gft_budget}
\end{figure}

% ============================================================
\subsection{Structural stability: GFT*}
\label{sec:gft_star}

The GFT is universal, but the prime sieve has a further property:
the partition ratio $\DKL/H$ converges.

\mDER
\begin{theorem}[Structural Stability, GFT* (conditional on PNT)]
\label{thm:gft_star}
For the prime gap distribution at primorial levels $m_k = 2 \cdot 3 \cdots p_k$,
the partition ratio $\DKL(P_{m_k})/H(P_{m_k})$ converges to a definite
limit as $k \to \infty$. The partition is structurally stable.

\begin{proof}
By the prime number theorem for arithmetic progressions (PNT-AP), the
distribution $P_{m_k}$ of primes modulo $m_k$ converges to the Euler
product measure $\mu_{m_k}(r) \propto \prod_{p|m_k}(1-1/p)^{-1}$
(for $\gcd(r, m_k) = 1$), which is uniform on reduced residues.
The KL divergence $\DKL(P_{m_k}\|\mu_{m_k})$ and entropy $H(P_{m_k})$
are continuous functions of $P_{m_k}$ and converge as $k \to \infty$.
Stability follows from the strict convexity of $\DKL$ and concavity of $H$.
The convergence rate is $O(1/\ln p_k)$ by PNT-AP; the ratio
$\DKL/H \to L$ where $L \approx s^2 = 1/4$.
\end{proof}
\end{theorem}

\begin{remark}
GFT* distinguishes the prime sieve from a random sieve. A random process
satisfies GFT-ID (universally), but the ratio $\DKL/H$ fluctuates without
converging. For primes, it converges to the Mertens constant.
\end{remark}

\paragraph{Complex extension.}
The action $S_{\mathrm{PT}} = -\sum\ln\sinp{p}$ extends to
$S_{\mathbb{C}} = -\sum\ln w_p$ with $\mathrm{Re}(S_{\mathbb{C}}) = S_{\mathrm{PT}}/2$,
where $w_p = (1 - e^{2i\theta_p})/2$ is the complex sieve variable.
See the companion article~\cite{companion_M4}.

% ============================================================
\subsection{Ruelle equivalence}
\label{sec:ruelle}

The transfer matrix $T_m$ connects GFT to Ruelle's thermodynamic
formalism~\cite{Ruelle1978}.

\mID
\begin{idbox}[Ruelle Equivalence]
\label{id:ruelle}
The Ruelle partition function of the sieve is:
\begin{equation}
  Z_{\mathrm{Ruelle}} = \mathrm{Tr}(T_m^N) = \sum_i \lambda_i^N,
  \label{eq:ruelle}
\end{equation}
where $\lambda_i$ are the eigenvalues of $T_m$. For a row-stochastic
matrix, $\lambda_0 = 1$ (Perron--Frobenius). The Ruelle free energy is:
\begin{equation}
  F_{\mathrm{Ruelle}} = -\ln \lambda_0 = 0.
  \label{eq:ruelle_free}
\end{equation}
The system is in equilibrium: $F_R = 0$.

\medskip\noindent\textbf{Status:} \CONDTAG\ (conditional on irreducibility of $T_m$
for all $m$; verified numerically for $m \leq 2310$, $N \leq 10^7$,
but not proved for all $m$).

\begin{proof}
$T_m$ restricted to the $\varphi(m)$ coprime residue classes is a
row-stochastic matrix (rows sum to~1, entries $\geq 0$).
It is empirically irreducible for all moduli tested ($m \leq 2310$,
primes up to $10^7$): every entry $(T_m)_{ab}$ with $\gcd(a,m) = \gcd(b,m) = 1$
is strictly positive.
(A proof of irreducibility for all sample sizes would require showing
that every coprime-class pair is realized by consecutive primes, which
is widely expected but not formally established.)
Assuming irreducibility, the Perron--Frobenius theorem guarantees a
unique maximal eigenvalue $\lambda_0 = 1$ with positive eigenvector
(stationary distribution $\pi$).
The trace formula $\mathrm{Tr}(T_m^N) = \sum_i
\lambda_i^N$ is the spectral decomposition of the trace.
$F_R = -\ln 1 = 0$ follows directly.
\end{proof}
\end{idbox}

\begin{remark}
The same partition function $Z = \mathrm{Tr}(T_m^N)$ appears in
Polyakov's transfer-matrix method~\cite{Polyakov1981}
(1D statistical mechanics) and in Regge calculus~\cite{Regge1959}. The triple coincidence $Z_{\mathrm{Feynman}} =
Z_{\mathrm{Ruelle}} = Z_{\mathrm{Polyakov}}$ is developed in
the companion article~\cite{companion_M5}.
\end{remark}

\begin{remark}[Topological pressure route: second proof of $F_R = 0$]
\label{rem:ruelle_pressure}
The Ruelle free energy $F_R = 0$ admits a second, structurally
independent derivation via the \emph{topological pressure} of the
Ruelle--Perron--Frobenius formalism~\cite{Ruelle1978}.

Define the potential $\varphi : \Sigma \to \mathbb{R}$ on the symbolic
shift $\Sigma$ induced by $T_m$ as $\varphi(\omega) = \ln T_m(\omega_0, \omega_1)$
(the log-transition weight).  The topological pressure is:
\[
  P(\varphi) = \lim_{n \to \infty} \frac{1}{n}
    \ln \sum_{\omega \in \Sigma_n} \exp\!\Bigl(\sum_{k=0}^{n-1}\varphi(\sigma^k\omega)\Bigr)
  = \ln \lambda_0(L_\varphi),
\]
where $L_\varphi$ is the Ruelle transfer operator and $\lambda_0(L_\varphi)$
its spectral radius.  For a row-stochastic matrix, $\lambda_0 = 1$
(Perron--Frobenius), hence $P(\varphi) = \ln 1 = 0$.
The variational principle then gives:
\[
  0 = P(\varphi) = \sup_\mu \bigl[h(\mu) + \int \varphi\,d\mu\bigr],
\]
where $h(\mu)$ is the measure-theoretic entropy.
Under the irreducibility assumption (empirically verified, see proof above),
the supremum is attained at the unique equilibrium state $\mu_{\rm eq} = \pi$
(the Perron eigenvector).
Note that $\pi$ is the measure of maximal entropy (MME) if and only if
$T_m$ is doubly stochastic; this holds at the sieve level for $m = 3$
(where $T_3$ is the antidiagonal permutation matrix) but need not hold
for the empirical prime transfer matrix at general~$m$.
The equality $F_R = -P(\varphi) = 0$
recovers Identity~\ref{id:ruelle} without using the trace formula
$\mathrm{Tr}(T^N)$---the argument depends only on the spectral
radius of $L_\varphi$, not on the matrix decomposition.

This provides a second independent route to $F_R = 0$: the first
uses the trace and Perron--Frobenius eigenvalue directly; the second
uses the variational principle of topological pressure.
Both routes are conditional on irreducibility of the restricted
matrix $T_m$ on coprime classes.
\end{remark}

% ============================================================
\subsection{Arrow of Time}
\label{sec:arrow}

Differentiating the GFT at fixed $m$ gives an irreversibility statement.

\mID
\begin{idbox}[Arrow of Time]
\label{id:arrow}
At fixed modulus $m$, along any path in distribution space satisfying
the GFT constraint:
\begin{equation}
  \frac{dH}{d\DKL} = -1.
  \label{eq:arrow}
\end{equation}

\begin{proof}
Differentiate $\DKL + H = \log_2 m$ (constant in $m$) with respect
to any parameter: $d\DKL + dH = 0$, hence $dH/d\DKL = -1$.

More precisely: as the prime distribution evolves toward the uniform
(as $N \to \infty$ by PNT-AP), $\DKL$ decreases while $H$ increases,
at equal and opposite rates along the GFT line.
\end{proof}
\end{idbox}

\begin{remark}
The Arrow of Time (dH/dDKL = -1) is exact: no inequality, no approximation.
The sieve's second law is an identity rather than an inequality.
\end{remark}

% ============================================================
\subsection{Bekenstein bound}
\label{sec:bekenstein}

\mTHM
\begin{theorem}[Bekenstein Bound]
\label{thm:bekenstein}
For any distribution $P$ on $m$ classes:
\begin{equation}
  \DKL(P\|U_m) \leq \log_2 m,
  \label{eq:bekenstein}
\end{equation}
with equality iff $P$ is a delta function on one class.

\begin{proof}
Since $H(P) \geq 0$ (entropy non-negative, with equality iff $P$ is
a point mass), the GFT gives $\DKL = \log_2 m - H \leq \log_2 m$.
Equality: $H = 0$ iff $P = \delta_r$ for some $r$.
\end{proof}
\end{theorem}

% ============================================================
\subsection{The derivation chain}
\label{sec:ch4_chain}

\[
  s = \tfrac{1}{2}
  \xrightarrow{\text{T1}} T_3
  \xrightarrow{\text{L0}} P_m \text{ geometric}
  \xrightarrow{\text{GFT}} \log_2 m = \DKL + H \xrightarrow{\text{Arrow}} dH/d\DKL = -1
\]

This chain is entirely within arithmetic. No physics enters.

% ============================================================
\subsection{Zero as cycle closure}
\label{sec:ch4_zero}

The GFT identity $\log_2 m = \DKL + H$ has a natural interpretation
through the concept of \textbf{cycle closure}.

\begin{remark}[Zero is not absence but phase cancellation]
\label{rem:zero_phase_gft}
In the sieve, the residue $n \equiv 0 \pmod{p}$ does not mean
``$n$ is absent at level~$p$'': it means the $p$-phase of~$n$
has \textbf{completed a full cycle} of $\mathbb{Z}/p\mathbb{Z}$
and returned to the origin.  A composite $n = k \cdot p$ has
traversed the circle $\mathbb{Z}/p\mathbb{Z}$ exactly $k$~times;
its holonomy has cancelled.

This gives a two-level reading of the sieve:
\begin{itemize}[nosep]
\item \textbf{Primes} $=$ perpetually open phases (no cycle closes
  below them).  They are irreducible precisely because they carry
  \emph{unclosed} holonomy.
\item \textbf{Composites} $=$ partially closed phases.
  $n = \prod p_i^{a_i}$ means the cycles of $\{p_i\}$ have each
  closed $a_i$ times.  The sieve removes them because a closed
  cycle carries no further open information.
\item \textbf{Zero} $=$ universal closure: every prime divides~$0$,
  so every cycle is simultaneously closed.  It is the
  \emph{informational vacuum}---not empty, but fully cancelled.
\item \textbf{One} $=$ no closure: no prime divides~$1$, so every
  cycle is fully open.  It is the \emph{identity operator}.
\end{itemize}

The GFT identity reflects this structure.  When $\DKL = 0$
(complete dissipation), all structure has been returned to
entropy: every informational cycle has closed.  When
$H = 0$ (zero entropy), all entropy has been converted to
structure: no cycle has closed, the state is maximally
persistent.  The conservation law $\DKL + H = \log_2 m$ says
that the total capacity is shared between open phases (structure)
and closed phases (entropy), with no loss.
\end{remark}

\begin{remark}[Fixed points as global cycle closures]
\label{rem:global_closure}
The fixed-point equation $\mustar = \sum_{\mathrm{active}} p$
(the companion article~\cite{companion_M3}) is a \textbf{global} cycle closure:
the sum of the active operators equals the scale that defines them.
The sieve ``recognises itself''---the structure it generates
(the active primes) sums to the scale that generates the structure.

Local closure: $n \equiv 0 \pmod{p}$ (one prime divides $n$).\\
Global closure: $\mustar = \sum p_{\rm active}$ (the sieve closes on itself).
\end{remark}

\begin{remark}[GFT as addition/multiplication conservation]
\label{rem:gft_add_mult}
The GFT identity $\log_2 m = \DKL + H$ can be read as a
\textbf{conservation law between the two operations}:
\begin{itemize}[nosep]
\item $\DKL$ (structure, persistence) measures
  \emph{multiplicative} information: how far the distribution
  deviates from uniform, quantified via log-ratios (logarithm
  converts $\times$ to $+$).
\item $H$ (entropy, disorder) measures \emph{additive}
  information: the sum $-\sum P \log P$ counts bits
  additively.
\item $\log_2 m$ (capacity) is the total: fixed, independent
  of how the information is partitioned.
\end{itemize}
The arrow identity $dH/d\DKL = -1$ says:
every unit of multiplicative structure gained costs exactly one
unit of additive entropy.  The two operations are in exact
balance.

This is the information-theoretic expression of the duality
described in \cref{rem:add_mult_duality}: the sieve
alternates $\times$ and $+$, and GFT conserves their sum.
\end{remark}

% ============================================================
\subsection{Summary and connections}
\label{sec:ch4_summary}

\begin{ledgerbox}
Hard core:
  GFT-ID ($\log_2 m = D_{\rm KL} + H$) -- exact algebraic tautology.
  Arrow ($dH/dD_{\rm KL} = -1$) -- exact differentiation.
  Bekenstein ($D_{\rm KL} \leq \log_2 m$) -- follows from $H \geq 0$.
  GFT* (structural stability) -- uses PNT-AP convergence.

Conditional:
  Ruelle ($F_R = 0$, $\mathrm{Tr}(T^N)$) -- conditional on irreducibility
  of $T_m$ for all $m$; verified numerically for $m \leq 2310$.

Bridge claim:
  Thermodynamic interpretation of $D_{\rm KL}$ and $H$ (the companion article~\cite{companion_M5}).
\end{ledgerbox}

\medskip\noindent\textbf{Forward connections.}
\begin{itemize}
\item The companion article~\cite{companion_M2} identifies $\DKL$ as the canonical divergence (Shore--Johnson).
\item The companion article~\cite{companion_M5} uses $Z_{\mathrm{Ruelle}} = \mathrm{Tr}(T_m^N)$ to
  construct the quantum partition function.
\item The companion article~\cite{companion_M5} promotes GFT to the first law of sieve thermodynamics.
\end{itemize}


% ============================================================
%  Conclusion
% ============================================================
\section{Conclusion}
\label{sec:conclusion}

This article has established the first four pillars of the Theory of
Persistence, proceeding from the single observation that the
Eratosthenes sieve forbids self-transitions modulo~3 among its
candidates.

\textbf{The sieve as a dynamical field} (Section~\ref{sec:sieve}).
The gap sequence $\{g_n\}$ is the unique degree of freedom.  Its
mod-$m$ projection defines a $\mathbb{Z}/m\mathbb{Z}$-valued gauge
connection, factoring via CRT into independent components at each
prime.  The forbidden-transition theorem (T1) forces the mod-$3$
transfer matrix to be purely antidiagonal, $T_3 = \mathrm{antidiag}(1,1)$,
with eigenvalues $\{+1,-1\}$ and stationary distribution $s = 1/2$.
Three independent routes---gap arithmetic, $\mathbb{Z}/6\mathbb{Z}$
involution, spectral constraint---establish $T_3$; the DPI
monotonicity orders the sieve levels by KL divergence.

\textbf{Uniqueness of Eratosthenes} (Section~\ref{sec:uniqueness}).
Four naturality theorems (N1--N4) prove that Eratosthenes is the
unique self-consistent, structurally minimal sieve with canonical
ordering; $p = 3$ is forced as the first cascade level and $s = 1/2$
as the first dynamical output.  Theorem T6 then provides three
methodologically distinct proofs of sieve uniqueness---from field
structure (T6a: the only proper ideal of $\mathbb{Z}/p\mathbb{Z}$ is
$\{0\}$), from four ring-congruence axioms C1--C4 (T6b: independence
verified by 4/4 counter-models), and from canonical information
geometry (T6c: $\DKL$ unique by Shore--Johnson, Fisher metric unique
by \v{C}encov).

\textbf{Conservation laws} (Section~\ref{sec:conservation}).
The maximum-entropy gap distribution is geometric (L0, by strict
concavity of entropy under a mean constraint).  The telescoping
conservation identity $\sum g_n = p_{N+1} - 2$ is the sieve's equation
of state.  The spectral conservation exponent $\alpha_{\rm cons} =
s^2 = 1/4$ (Theorem T2) controls the exponential convergence rate,
established both by direct spectral computation on $T_{30}$ and by a
CRT trace formula.  The vertex--edge bifurcation at $\mustar = 15$
produces exactly two canonical parameters, $\qrel = 13/15$ and
$\qtherm = e^{-1/15}$, derived independently via three routes
(optimisation, exponential-family duality, partition function).

\textbf{The Gap Fundamental Theorem} (Section~\ref{sec:gft}).
The exact algebraic identity $\log_2 m = \DKL + H$ decomposes the
total information capacity into persistent structure and irreducible
noise.  The Arrow of Time $dH/d\DKL = -1$ is an exact identity, not
an inequality.  The Ruelle identification connects the GFT to
thermodynamic formalism: the sieve's transfer matrix has zero free
energy ($F_R = 0$), and its stationary distribution is the unique
equilibrium measure.  The Bekenstein bound $\DKL \leq \log_2 m$
follows as a corollary.

\paragraph{The toric spiral.}
The number line is a 1D projection of a richer structure.
The CRT decomposition maps the integers onto the torus
$T^3 = \mathbb{Z}/3\mathbb{Z} \times \mathbb{Z}/5\mathbb{Z} \times
\mathbb{Z}/7\mathbb{Z}$, where the prime gap sequence traces a
quasi-periodic spiral.  The convergence $\alpha_k \to 1/2$
(Article~III~\cite{companion_M3}) is the geometric tightening of this
spiral toward the limit circle.  The three spatial dimensions of the
Bianchi~I metric (Article~V~\cite{companion_M5}) correspond to the
three circles of this torus.  This interpretation unifies the
algebraic (CRT), dynamical (convergence), and geometric (metric)
aspects of the theory in a single visual picture.

\medskip
\noindent\textbf{Open problems.}
\begin{enumerate}[nosep]
\item \emph{Convergence rate of $\alpha_k$.}  The value $s = 1/2$ is
  the unique stationary distribution of $T_3$ (a theorem of linear
  algebra).  The rate at which the empirical frequency $\alpha_k$
  converges to $s$ is quantified in Article~III~\cite{companion_M3}
  via spectral annihilation and Mertens compactness.
\item \emph{Irreducibility of $T_m$.}  The Ruelle identification and
  GFT* both assume irreducibility of the transfer matrix on coprime
  classes.  A proof for all $m$ remains open.
\item \emph{Fixed point $\mustar = 15$.}  The self-consistency scan
  identifies $\mustar = 15$ as the unique cascade fixed point; an
  analytic proof is developed in the companion article~\cite{companion_M3}.
\item \emph{Physical interpretation.}  All results in this article are
  purely arithmetic.  Whether the vertex--edge bifurcation, the
  conservation laws, and the GFT identity have physical content is
  the subject of the companion article~\cite{companion_M5}.
\end{enumerate}

\medskip
\noindent\textbf{What comes next.}
The companion article~\cite{companion_M2} develops the canonical information geometry
(Fisher metric, holonomy angles $\sin^2\theta_p$) and the internal
angle scale.  The companion article~\cite{companion_M3} establishes the exact recurrences,
spectral annihilation, and the fixed-point theorem $\mustar = 15$.
Together, the five articles provide the complete mathematical
foundation.  The bridge to physical observables is
developed in Article~V~\cite{companion_M5}.

\begin{table}[htbp]
\centering
\caption{Dependency and epistemic status of main results.}
\small
\begin{tabular}{@{}llll@{}}
\toprule
\textbf{Result} & \textbf{Depends on} & \textbf{Status} & \textbf{Notes} \\
\midrule
T1 (Forbidden Trans.) & None & \THMTAG & Unconditional \\
T3 (Antidiagonal) & T1 & \THMTAG & Unconditional \\
$s = 1/2$ & T3 (Perron--Frobenius) & \THMTAG & Unique stationary distribution \\
N1--N2 (Uniqueness) & FTA & \THMTAG & Classical \\
N3--N4 (Minimality) & N1--N2 & \THMTAG & Novel \\
T6 (Sieve uniqueness) & N1--N4 & \THMTAG & Three proofs \\
L0 (Max entropy) & Lagrange, Jaynes & \THMTAG & Standard technique \\
T2 (Spectral cons.) & $T_{30}$ spectrum & \VALTAG & Numerical; alg.\ proof open \\
Bifurcation & L0 + $\mustar$ (M3) & \CONDTAG & Conditional on M3 \\
GFT identity & Shannon decomp. & \IDTAG & Valid for any distribution \\
Ruelle equivalence & Irreducibility & \CONDTAG & Open for all $m$ \\
\bottomrule
\end{tabular}
\end{table}


\bibliographystyle{plainnat}
\bibliography{references}

\end{document}
